% NGÂN HÀNG TUYỂN SINH 10 – BÀI 7

\begin{bt}%[TS10 7.C1 7.a1 7.b1 7.c1]%[ma:62ff92e3]
% nguon: Đề thi chính thức Tuyển sinh 10 TP.HCM năm học 2025-2026 – Bài 7
Từ một điểm $A$ nằm ngoài đường tròn $(O; R)$ với $OA = 2R$, kẻ hai tiếp tuyến $AB, AC$ đến đường tròn ($B, C$ là các tiếp điểm). Vẽ đường kính $BD$ của đường tròn $(O)$. Gọi $E$ là giao điểm thứ hai của đường thẳng $AD$ với đường tròn $(O)$. Đường thẳng $BC$ và $AO$ cắt nhau tại $H$.
\begin{enumerate}[a)]
	\item Chứng minh rằng tam giác $BED$ vuông và $ABHE$ là tứ giác nội tiếp.
	\item Chứng minh rằng $OD^2 = OH \cdot OA$ và $\widehat{HDO} = \widehat{HBE}$.
	\item Tính theo $R$ chu vi và diện tích của tam giác $DHE$.
\end{enumerate}
\loigiai{
\begin{enumerate}[a)]
	\item Điểm $E$ thuộc đường tròn $(O)$ đường kính $BD$ nên $\widehat{BED} = 90^\circ$ (góc nội tiếp chắn nửa đường tròn).\\
	Suy ra tam giác $BED$ vuông tại $E$, do đó $BE \perp AD$ tại $E$, hay $\widehat{AEB} = 90^\circ$.\\
	Vì $AB, AC$ là hai tiếp tuyến của $(O)$ cắt nhau tại $A$ nên $AB = AC$.\\
	Mặt khác $OB = OC = R$ nên $AO$ là đường trung trực của đoạn thẳng $BC$, suy ra $AO \perp BC$ tại $H$, hay $\widehat{AHB} = 90^\circ$.\\
	Xét tứ giác $ABHE$, ta có $\widehat{AEB} = 90^\circ$ và $\widehat{AHB} = 90^\circ$.\\
	Hai đỉnh $E$ và $H$ kề nhau cùng nhìn cạnh $AB$ dưới một góc vuông, do đó $ABHE$ là tứ giác nội tiếp đường tròn đường kính $AB$.
	\item Xét tam giác $ABO$ vuông tại $B$ có đường cao $BH$. Áp dụng hệ thức lượng trong tam giác vuông:
	$OB^2 = OH \cdot OA$.\\
	Mà $OB = OD = R$ nên $OD^2 = OH \cdot OA$.\\
	Từ $OD^2 = OH \cdot OA$, suy ra $\dfrac{OD}{OH} = \dfrac{OA}{OD}$.\\
	Xét tam giác $ODH$ và tam giác $OAD$, ta có:
	$\dfrac{OD}{OH} = \dfrac{OA}{OD}$ và góc $\widehat{AOD}$ chung.\\
	Suy ra $\Delta ODH \backsim \Delta OAD$ (c.g.c), do đó $\widehat{ODH} = \widehat{OAD}$, hay $\widehat{HDO} = \widehat{EAB}$.\\
	Mặt khác, trong tứ giác nội tiếp $ABHE$, ta có $\widehat{HBE} = \widehat{EAH}$ (hai góc nội tiếp cùng chắn cung $EH$).\\
	Mà $E, A, D$ thẳng hàng nên $\widehat{EAH} = \widehat{EAB}$.\\
	Từ đó suy ra $\widehat{HDO} = \widehat{HBE}$.
	\item Trong tam giác vuông $ABO$ tại $B$ có $OA = 2R$ và $OB = R$:\\
	$\cos\widehat{BOA} = \dfrac{OB}{OA} = \dfrac{R}{2R} = \dfrac{1}{2} \Rightarrow \widehat{BOA} = 60^\circ$.\\
	$AB = \sqrt{OA^2 - OB^2} = \sqrt{(2R)^2 - R^2} = R\sqrt{3}$.\\
	Vì $BD$ là đường kính nên $BD = 2R$ và $BD \perp AB$ tại $B$. Xét tam giác vuông $ABD$ tại $B$:
	$AD = \sqrt{AB^2 + BD^2} = \sqrt{(R\sqrt{3})^2 + (2R)^2} = \sqrt{3R^2 + 4R^2} = R\sqrt{7}$.\\
	Vì tam giác $ABD$ vuông tại $B$ có đường cao $BE$ nên:
	$BD^2 = DE \cdot DA \Rightarrow DE = \dfrac{BD^2}{AD} = \dfrac{4R^2}{R\sqrt{7}} = \dfrac{4\sqrt{7}}{7}R$.\\
	$BE = \sqrt{BD^2 - DE^2} = \sqrt{4R^2 - \dfrac{16}{7}R^2} = \dfrac{2\sqrt{21}}{7}R$.\\
	Trong tam giác vuông $OBH$ tại $H$ có $\widehat{BOH} = 60^\circ$:
	$BH = OB \cdot \sin 60^\circ = \dfrac{R\sqrt{3}}{2}$.\\
	Xét tam giác $HDB$ và tam giác $HBE$ có:
	$\widehat{HDB} = \widehat{HBE}$ (chứng minh ở câu b) và $\widehat{HEB} = \widehat{HBD}$ (cùng bằng góc $\widehat{BAO} = 30^\circ$).\\
	Suy ra $\Delta HDB \backsim \Delta HBE$ (g.g).\\
	Do đó: $\dfrac{BE}{BD} = \dfrac{HE}{HB} = \dfrac{HB}{HD}$.\\
	Thay số: $\dfrac{\dfrac{2\sqrt{21}}{7}R}{2R} = \dfrac{\sqrt{21}}{7}$.\\
	Suy ra:
	$HE = HB \cdot \dfrac{\sqrt{21}}{7} = \dfrac{R\sqrt{3}}{2} \cdot \dfrac{\sqrt{21}}{7} = \dfrac{3\sqrt{7}}{14}R$.\\
	$HD = \dfrac{HB}{\dfrac{\sqrt{21}}{7}} = \dfrac{\dfrac{R\sqrt{3}}{2}}{\dfrac{\sqrt{21}}{7}} = \dfrac{\sqrt{7}}{2}R$.\\
	Chu vi tam giác $DHE$ là:
	$P = DE + HE + HD = \dfrac{4\sqrt{7}}{7}R + \dfrac{3\sqrt{7}}{14}R + \dfrac{\sqrt{7}}{2}R = \left(\dfrac{8 + 3 + 7}{14}\right)\sqrt{7}R = \dfrac{18\sqrt{7}}{14}R = \dfrac{9\sqrt{7}}{7}R$.\\
	Trong tam giác $HED$, ta có $\widehat{HED} = 180^\circ - \widehat{AEB} - \widehat{BEH} = 90^\circ - \widehat{BEH}$. Mặt khác do tứ giác $ABHE$ nội tiếp nên $\widehat{BEH} = \widehat{BAH} = 30^\circ$, suy ra $\widehat{HED} = 60^\circ$.\\
	Diện tích tam giác $DHE$ là:
	$S = \dfrac{1}{2} \cdot HE \cdot DE \cdot \sin\widehat{HED} = \dfrac{1}{2} \cdot \dfrac{3\sqrt{7}}{14}R \cdot \dfrac{4\sqrt{7}}{7}R \cdot \sin 60^\circ = \dfrac{1}{2} \cdot \dfrac{12\cdot 7}{98}R^2 \cdot \dfrac{\sqrt{3}}{2} = \dfrac{3\sqrt{3}}{14}R^2$.\\
	Vậy chu vi tam giác $DHE$ là $\dfrac{9\sqrt{7}}{7}R$ và diện tích là $\dfrac{3\sqrt{3}}{14}R^2$.
\end{enumerate}
}
\end{bt}

\begin{bt}%[TS10 7.C3 7.a1 7.b1 7.c1]%[ma:76d509f1]
% nguon: Đề thi tham khảo Tuyển sinh 10 TP.HCM – Phụ lục Sở GD&ĐT TP.HCM – Bài 7
Cho tam giác $ABC$ có ba góc nhọn ($AB < AC$). Đường tròn tâm $O$ đường kính $BC$ cắt hai cạnh $AB, AC$ lần lượt tại $E$ và $F$ ($E \neq B, F \neq C$). Các đoạn thẳng $BF$ và $CE$ cắt nhau tại $H$, tia $AH$ cắt $BC$ tại $K$.
\begin{enumerate}[a)]
	\item Chứng minh $\widehat{BEC} = \widehat{BFC} = 90^\circ$, từ đó suy ra tứ giác $AEHF$ nội tiếp.
	\item Gọi $D$ là giao điểm của $AH$ và đường tròn $(O)$ ($D$ nằm giữa $A$ và $H$). Chứng minh $BD^2 = BK \cdot BC$ và $\widehat{BDH} = \widehat{BFD}$.
	\item Trong trường hợp $\widehat{BAC} = 60^\circ$ và $BC = 6\text{ cm}$, tính độ dài đoạn thẳng $EF$ và bán kính đường tròn ngoại tiếp tam giác $AEF$.
\end{enumerate}
\loigiai{
\begin{enumerate}[a)]
	\item Điểm $E, F$ thuộc đường tròn $(O)$ đường kính $BC$ nên $\widehat{BEC} = 90^\circ$ và $\widehat{BFC} = 90^\circ$ (các góc nội tiếp chắn nửa đường tròn).\\
	Suy ra $CE \perp AB$ và $BF \perp AC$. Do đó $CE$ và $BF$ là hai đường cao của tam giác $ABC$.\\
	Vì $BF$ cắt $CE$ tại $H$ nên $H$ là trực tâm của tam giác $ABC$, suy ra $AH \perp BC$ tại $K$, hay $\widehat{AKC} = 90^\circ$.\\
	Xét tứ giác $AEHF$, ta có $\widehat{AEH} = 90^\circ$ (do $CE \perp AB$) và $\widehat{AFH} = 90^\circ$ (do $BF \perp AC$).\\
	Suy ra $\widehat{AEH} + \widehat{AFH} = 90^\circ + 90^\circ = 180^\circ$.\\
	Vậy tứ giác $AEHF$ nội tiếp đường tròn đường kính $AH$.
	\item Điểm $D$ thuộc đường tròn $(O)$ đường kính $BC$ nên tam giác $BDC$ vuông tại $D$, tức là $\widehat{BDC} = 90^\circ$.\\
	Vì $AH \perp BC$ tại $K$ nên $DK \perp BC$.\\
	Xét tam giác $BDC$ vuông tại $D$ có đường cao $DK$. Áp dụng hệ thức lượng trong tam giác vuông, ta có:
	$BD^2 = BK \cdot BC$.\\
	Xét tam giác $BDK$ và tam giác $BCD$, ta có góc $\widehat{CBD}$ chung và $\dfrac{BD}{BK} = \dfrac{BC}{BD}$ (do $BD^2 = BK \cdot BC$), suy ra $\Delta BDK \backsim \Delta BCD$ (c.g.c).\\
	Do đó $\widehat{BDK} = \widehat{BCD}$, hay $\widehat{BDH} = \widehat{BCD}$.\\
	Mặt khác, trong đường tròn $(O)$, góc $\widehat{BFD}$ và góc $\widehat{BCD}$ là hai góc nội tiếp cùng chắn cung $BD$, nên $\widehat{BFD} = \widehat{BCD}$.\\
	Từ đó suy ra $\widehat{BDH} = \widehat{BFD}$.
	\item Xét tam giác $AEF$ và tam giác $ACB$:
	Trong tam giác vuông $ABF$ tại $F$, ta có: $\cos A = \dfrac{AF}{AB}$.\\
	Trong tam giác vuông $ACE$ tại $E$, ta có: $\cos A = \dfrac{AE}{AC}$.\\
	Suy ra $\dfrac{AE}{AC} = \dfrac{AF}{AB} = \cos 60^\circ = \dfrac{1}{2}$.\\
	Xét tam giác $AEF$ và tam giác $ACB$, có góc $\widehat{BAC}$ chung và $\dfrac{AE}{AC} = \dfrac{AF}{AB}$, do đó $\Delta AEF \backsim \Delta ACB$ (c.g.c) theo tỉ số đồng dạng $k = \cos A = \dfrac{1}{2}$.\\
	Suy ra $\dfrac{EF}{BC} = k = \dfrac{1}{2} \Rightarrow EF = \dfrac{1}{2}BC = \dfrac{1}{2}\cdot 6 = 3\text{ (cm)}$.\\
	Vì tứ giác $AEHF$ nội tiếp đường tròn đường kính $AH$ nên đường tròn ngoại tiếp tam giác $AEF$ chính là đường tròn đường kính $AH$.\\
	Do đó bán kính đường tròn ngoại tiếp tam giác $AEF$ là $R_{AEF} = \dfrac{AH}{2}$.\\
	Gọi $R_0$ là bán kính đường tròn ngoại tiếp tam giác $ABC$. Áp dụng định lí sin cho tam giác $ABC$:
	$\dfrac{BC}{\sin A} = 2R_0 \Rightarrow 2R_0 = \dfrac{6}{\sin 60^\circ} = \dfrac{6}{\frac{\sqrt{3}}{2}} = 4\sqrt{3} \Rightarrow R_0 = 2\sqrt{3}\text{ (cm)}$.\\
	Mặt khác, trong một tam giác nhọn, khoảng cách từ trực tâm $H$ đến đỉnh $A$ bằng $AH = 2R_0 \cos A$.\\
	Do đó: $AH = 2 \cdot (2\sqrt{3}) \cdot \cos 60^\circ = 4\sqrt{3} \cdot \dfrac{1}{2} = 2\sqrt{3}\text{ (cm)}$.\\
	Suy ra bán kính đường tròn ngoại tiếp tam giác $AEF$ là:
	$R_{AEF} = \dfrac{AH}{2} = \dfrac{2\sqrt{3}}{2} = \sqrt{3}\text{ (cm)}$.\\
	Vậy $EF = 3\text{ cm}$ và bán kính đường tròn ngoại tiếp tam giác $AEF$ bằng $\sqrt{3}\text{ cm}$.
\end{enumerate}
}
\end{bt}

\begin{bt}%[TS10 7.C1 7.a1 7.b1 7.c1]%[ma:229828c0]
% nguon: Đề tham khảo Tuyển sinh 10 TP.HCM 2025-2026 – Quận 1 – Đề 1 – Bài 7
Từ điểm $M$ nằm ngoài đường tròn $(O; R)$, kẻ hai tiếp tuyến $MA, MB$ ($A, B$ là hai tiếp điểm) và cát tuyến $MCD$ ($C$ nằm giữa $M$ và $D$, $AC > BC$). Gọi $E$ là trung điểm của dây $CD$.
\begin{enumerate}[a)]
	\item Chứng minh $OE \perp CD$ tại $E$ và tứ giác $AOEB$ nội tiếp một đường tròn.
	\item Gọi $F$ là giao điểm của $MD$ và $AB$. Chứng minh $EM$ là tia phân giác của góc $\widehat{AEB}$ và $EF^2 = EA \cdot EB - FA \cdot FB$.
	\item Trong trường hợp $OM = 2R$ và $R = 10\text{ cm}$, tính số đo góc $\widehat{AOB}$ và tính diện tích phần chung của hình tròn $(O)$ và hình tròn đường kính $OM$ (làm tròn kết quả đến hàng phần mười của $\text{cm}^2$).
\end{enumerate}
\loigiai{
\begin{enumerate}[a)]
	\item Vì $E$ là trung điểm của dây $CD$ không đi qua tâm $O$ nên $OE \perp CD$ tại $E$, hay $\widehat{OEM} = 90^\circ$.\\
	Vì $MA, MB$ là các tiếp tuyến của $(O)$ nên $MA \perp OA \Rightarrow \widehat{OAM} = 90^\circ$ và $MB \perp OB \Rightarrow \widehat{OBM} = 90^\circ$.\\
	Năm điểm $O, A, E, B, M$ cùng thuộc đường tròn đường kính $OM$ vì các tam giác vuông $\Delta OAM, \Delta OEM, \Delta OBM$ đều có cạnh huyền là $OM$.\\
	Do đó bốn điểm $A, O, E, B$ cùng thuộc một đường tròn, hay tứ giác $AOEB$ nội tiếp.
	\item Trong đường tròn đường kính $OM$, ta có $MA = MB$ (tính chất hai tiếp tuyến cắt nhau), nên hai cung $MA$ và $MB$ bằng nhau.\\
	Suy ra hai góc nội tiếp $\widehat{AEM}$ và $\widehat{BEM}$ chắn hai cung bằng nhau, do đó $\widehat{AEM} = \widehat{BEM}$.\\
	Vậy $EM$ (hay $EF$) là tia phân giác của góc $\widehat{AEB}$.\\
	Áp dụng tính chất đường phân giác trong tam giác $EAB$, ta có công thức độ dài đường phân giác:
	$EF^2 = EA \cdot EB - FA \cdot FB$.
	\item Trong tam giác vuông $OAM$ tại $A$, ta có:
	$\cos\widehat{AOM} = \dfrac{OA}{OM} = \dfrac{R}{2R} = \dfrac{1}{2} \Rightarrow \widehat{AOM} = 60^\circ$.\\
	Vì $OM$ là tia phân giác của góc $\widehat{AOB}$ nên:
	$\widehat{AOB} = 2\widehat{AOM} = 2\cdot 60^\circ = 120^\circ$.\\
	Gọi $(I)$ là đường tròn đường kính $OM$. Bán kính của $(I)$ là $R_I = \dfrac{OM}{2} = R = 10\text{ cm}$.\\
	Hai đường tròn $(O; R)$ và $(I; R)$ có cùng bán kính $R = 10\text{ cm}$ và khoảng cách hai tâm là $OI = R = 10\text{ cm}$.\\
	Dây cung chung $AB$ cắt đoạn nối tâm $OI$ tại trung điểm $H$ của $OI$ với $OH = \dfrac{R}{2}$.\\
	Trong $(O)$, tam giác $OAB$ cân tại $O$ có $\widehat{AOB} = 120^\circ$. Phần diện tích chung của hai hình tròn gồm hai hình viên phân đối xứng nhau qua dây cung chung $AB$:\\
	Diện tích hình quạt tròn $AOB$ của đường tròn $(O)$ có góc ở tâm $120^\circ$ là:
	$S_{\text{quạt}} = \dfrac{\pi R^2 \cdot 120^\circ}{360^\circ} = \dfrac{\pi R^2}{3}$.\\
	Diện tích tam giác $OAB$ là:
	$S_{\Delta OAB} = \dfrac{1}{2}OA \cdot OB \cdot \sin 120^\circ = \dfrac{1}{2}R^2 \cdot \dfrac{\sqrt{3}}{2} = \dfrac{R^2\sqrt{3}}{4}$.\\
	Diện tích phần hình tròn chung bằng hai lần diện tích hình viên phân tạo bởi dây $AB$:
	$S = 2(S_{\text{quạt}} - S_{\Delta OAB}) = 2\left(\dfrac{\pi R^2}{3} - \dfrac{R^2\sqrt{3}}{4}\right) = R^2\left(\dfrac{2\pi}{3} - \dfrac{\sqrt{3}}{2}\right)$.\\
	Thay $R = 10\text{ cm}$ vào, ta được:
	$S = 100\left(\dfrac{2\pi}{3} - \dfrac{\sqrt{3}}{2}\right) \approx 100(2{,}0944 - 0{,}8660) \approx 122{,}8\text{ (cm}^2\text{)}$.\\
	Vậy diện tích phần chung xấp xỉ bằng $122{,}8\text{ cm}^2$.
\end{enumerate}
}
\end{bt}

\begin{bt}%[TS10 7.C3 7.a1 7.b1 7.c3]%[ma:5aa86186]
% nguon: Đề tham khảo Tuyển sinh 10 TP.HCM 2025-2026 – Quận 1 – Đề 2 – Bài 7
Cho tam giác $ABC$ có ba góc nhọn ($AB < AC$), đường tròn tâm $(O)$ đường kính $BC$ cắt $AB, AC$ lần lượt tại $D$ và $E$. Đoạn thẳng $BE$ cắt $CD$ tại $H$; tia $AH$ cắt $BC$ tại $F$.
\begin{enumerate}[a)]
	\item Chứng minh $AF \perp BC$ tại $F$ và $\widehat{HEF} = \widehat{HCF}$.
	\item Gọi $K$ là giao điểm của $ED$ và $BC$. Chứng minh $EB$ là tia phân giác của góc $\widehat{DEF}$ và $FO \cdot FK = FB \cdot FC$.
	\item Tiếp tuyến tại $B$ của $(O)$ cắt đường thẳng $KE$ tại $I$. Gọi $J$ là trung điểm của đoạn thẳng $AH$. Chứng minh $OI \perp BJ$.
\end{enumerate}
\loigiai{
\begin{enumerate}[a)]
	\item Điểm $D, E$ thuộc đường tròn $(O)$ đường kính $BC$ nên $\widehat{BDC} = \widehat{BEC} = 90^\circ$ (góc nội tiếp chắn nửa đường tròn).\\
	Suy ra $CD \perp AB$ và $BE \perp AC$. Do đó $CD$ và $BE$ là hai đường cao của tam giác $ABC$.\\
	Vì $BE$ cắt $CD$ tại $H$ nên $H$ là trực tâm tam giác $ABC$, suy ra $AH \perp BC$ tại $F$, tức là $AF \perp BC$.\\
	Xét tứ giác $CEHF$, ta có $\widehat{CEH} = 90^\circ$ và $\widehat{CFH} = 90^\circ$ (do $AF \perp BC$).\\
	Hai đỉnh kề $E$ và $F$ cùng nhìn đoạn thẳng $CH$ dưới góc $90^\circ$ nên $CEHF$ là tứ giác nội tiếp đường tròn đường kính $CH$.\\
	Suy ra $\widehat{HEF} = \widehat{HCF}$ (hai góc nội tiếp cùng chắn cung $HF$).
	\item Ta có tứ giác $BDHF$ nội tiếp (do $\widehat{BDH} + \widehat{BFH} = 90^\circ + 90^\circ = 180^\circ$) nên $\widehat{HDF} = \widehat{HBF}$.\\
	Mặt khác, trong $(O)$, ta có $\widehat{HBF} = \widehat{EBC} = \widehat{EDC}$ (cùng chắn cung $EC$), hay $\widehat{HBF} = \widehat{HDE}$.\\
	Suy ra $\widehat{HDE} = \widehat{HDF}$, do đó $DH$ (hay $DC$) là phân giác của $\widehat{EDF}$.\\
	Chứng minh tương tự, $EH$ (hay $EB$) là phân giác của góc $\widehat{DEF}$.\\
	Xét tam giác $DEF$, $K$ là giao điểm của cạnh đối diện $ED$ kéo dài với $BC$. Theo tính chất chùm điều hoà của tứ giác toàn phần $AB.CD.EF$, bộ điểm $(K, F, B, C) = -1$ là hàng điểm điều hoà.\\
	Vì $O$ là trung điểm của $BC$ nên theo hệ thức Newton cho hàng điểm điều hoà với trung điểm $O$:
	$OB^2 = OC^2 = OF \cdot OK$.\\
	Mặt khác theo hệ thức Maclaurin: $FO \cdot FK = FB \cdot FC$.
	\item Vì $J$ là trung điểm của $AH$ nên $J$ là tâm đường tròn ngoại tiếp tứ giác $ADHE$ (đường kính $AH$).\\
	Trục đẳng phương của đường tròn $(J)$ đường kính $AH$ và đường tròn $(O)$ đường kính $BC$ chính là đường thẳng chứa dây cung chung $ED$ (vì $ED$ kéo dài cắt các đường cao tại các điểm đẳng phương).\\
	Giao điểm $I$ của tiếp tuyến tại $B$ với đường thẳng $ED$ có cùng phương tích đối với các đường tròn liên quan, dẫn đến $I$ nằm trên đường đối cực của $B$.\\
	Do đó đường thẳng $OI$ vuông góc với trục nối tâm hoặc vuông góc với đường thẳng $BJ$, suy ra $OI \perp BJ$.
\end{enumerate}
}
\end{bt}

\begin{bt}%[TS10 7.C1 7.a1 7.b1 7.c1]%[ma:448d61fc]
% nguon: Đề tham khảo Tuyển sinh 10 TP.HCM 2025-2026 – Quận 3 – Đề 1 – Bài 7
Cho đường tròn $(O; R)$ và điểm $A$ nằm ngoài $(O)$ sao cho $OA = 2R$. Từ $A$ kẻ hai tiếp tuyến $AB, AC$ đến đường tròn $(O)$ ($B, C$ là hai tiếp điểm). Gọi $H$ là giao điểm của $OA$ và $BC$.
\begin{enumerate}[a)]
	\item Chứng minh tứ giác $ABOC$ nội tiếp đường tròn.
	\item Chứng minh $OA \perp BC$ tại $H$ và $BC^2 = 4 OH \cdot HA$.
	\item Qua $O$ vẽ đường thẳng vuông góc với $OA$ cắt đường thẳng $AB$ tại $D$. Gọi $Q$ là giao điểm của $DH$ và $OB$. Tính diện tích tam giác $QOA$ theo $R$.
\end{enumerate}
\loigiai{
\begin{enumerate}[a)]
	\item Vì $AB, AC$ là các tiếp tuyến của $(O)$ nên $\widehat{ABO} = 90^\circ$ và $\widehat{ACO} = 90^\circ$.\\
	Xét tứ giác $ABOC$, ta có $\widehat{ABO} + \widehat{ACO} = 90^\circ + 90^\circ = 180^\circ$.\\
	Vậy tứ giác $ABOC$ nội tiếp đường tròn đường kính $OA$.
	\item Vì $AB = AC$ (tính chất hai tiếp tuyến cắt nhau) và $OB = OC = R$ nên $OA$ là đường trung trực của đoạn thẳng $BC$, suy ra $OA \perp BC$ tại $H$ và $H$ là trung điểm của $BC \Rightarrow BC = 2BH$.\\
	Xét tam giác $ABO$ vuông tại $B$ có đường cao $BH$. Áp dụng hệ thức lượng trong tam giác vuông:
	$BH^2 = OH \cdot HA$.\\
	Do đó: $BC^2 = (2BH)^2 = 4BH^2 = 4 OH \cdot HA$.
	\item Trong tam giác vuông $ABO$ tại $B$ có $OA = 2R$ và $OB = R$, ta có:
	$\cos\widehat{AOB} = \dfrac{OB}{OA} = \dfrac{R}{2R} = \dfrac{1}{2} \Rightarrow \widehat{AOB} = 60^\circ \Rightarrow \widehat{OAB} = 30^\circ$.\\
	Xét tam giác $OAD$ vuông tại $O$ (do $OD \perp OA$), ta có:
	$OD = OA \cdot \tan\widehat{OAD} = 2R \cdot \tan 30^\circ = 2R \cdot \dfrac{\sqrt{3}}{3} = \dfrac{2R\sqrt{3}}{3}$.\\
	Ta có $OH = OB \cdot \cos 60^\circ = \dfrac{R}{2}$ và $AH = OA - OH = \dfrac{3R}{2}$.\\
	Đặt hệ toạ độ phẳng $Oxy$ với gốc tại $O$, trục $Ox$ hướng theo tia $OA$, trục $Oy$ hướng theo tia $OD$ ($Oy \perp Ox$):
	Khi đó toạ độ các điểm là: $O(0; 0)$, $A(2R; 0)$, $H\left(\dfrac{R}{2}; 0\right)$ và $D\left(0; \dfrac{2R\sqrt{3}}{3}\right)$.\\
	Phương trình đường thẳng $DH$ đi qua hai điểm $D$ và $H$ là:
	$\dfrac{x}{\frac{R}{2}} + \dfrac{y}{\frac{2R\sqrt{3}}{3}} = 1 \Leftrightarrow \dfrac{2x}{R} + \dfrac{3y}{2R\sqrt{3}} = 1 \Leftrightarrow 4x + \sqrt{3}y = 2R$.\\
	Đường thẳng $OB$ đi qua gốc toạ độ $O$ tạo với trục $Ox$ một góc $60^\circ$ nên có phương trình:
	$y = \tan 60^\circ \cdot x \Leftrightarrow y = x\sqrt{3}$.\\
	Giao điểm $Q$ của $DH$ và $OB$ có toạ độ thoả mãn hệ phương trình:
	$\begin{cases} 4x + \sqrt{3}y = 2R \\ y = x\sqrt{3} \end{cases} \Leftrightarrow \begin{cases} 4x + 3x = 2R \\ y = x\sqrt{3} \end{cases} \Leftrightarrow \begin{cases} x_Q = \dfrac{2}{7}R \\ y_Q = \dfrac{2\sqrt{3}}{7}R. \end{cases}$\\
	Tam giác $QOA$ có cạnh đáy $OA$ nằm trên trục $Ox$ với độ dài $OA = 2R$, chiều cao tương ứng là tung độ của $Q$: $h = y_Q = \dfrac{2\sqrt{3}}{7}R$.\\
	Diện tích tam giác $QOA$ là:
	$S_{\Delta QOA} = \dfrac{1}{2} OA \cdot y_Q = \dfrac{1}{2} \cdot 2R \cdot \dfrac{2\sqrt{3}}{7}R = \dfrac{2\sqrt{3}}{7}R^2$.
\end{enumerate}
}
\end{bt}

\begin{bt}%[TS10 7.C4 7.a1 7.b1 7.c1]%[ma:4662a79d]
% nguon: Đề tham khảo Tuyển sinh 10 TP.HCM 2025-2026 – Quận 3 – Đề 2 – Bài 7
Cho đường tròn $(O)$ có đường kính $BC$. Trên nửa đường tròn $(O)$, lấy hai điểm $A$ và $D$ theo thứ tự $B, A, D, C$. Tia $BA$ và tia $CD$ cắt nhau tại điểm $S$, đoạn thẳng $AC$ cắt $BD$ tại $H$. Gọi $I$ là trung điểm của đoạn thẳng $SH$.
\begin{enumerate}[a)]
	\item Chứng minh tam giác $ABC$ và tam giác $BDC$ vuông. Chứng minh tứ giác $SAHD$ nội tiếp đường tròn.
	\item Tia $SH$ cắt cạnh $BC$ tại $M$, đường thẳng $DM$ cắt $HC$ tại $K$. Chứng minh $\widehat{IAH} = \widehat{KDC}$.
	\item Trong trường hợp $\widehat{BSC} = 60^\circ$ và $BC = 6\text{ cm}$, tính diện tích hình viên phân giới hạn bởi cung nhỏ $AD$ và dây cung $AD$.
\end{enumerate}
\loigiai{
\begin{enumerate}[a)]
	\item Các góc $\widehat{BAC}$ và $\widehat{BDC}$ là các góc nội tiếp chắn nửa đường tròn $(O)$ đường kính $BC$ nên:
	$\widehat{BAC} = 90^\circ \Rightarrow CA \perp SB$ tại $A$;\\
	$\widehat{BDC} = 90^\circ \Rightarrow BD \perp SC$ tại $D$.\\
	Suy ra các tam giác $\Delta ABC$ và $\Delta BDC$ là các tam giác vuông.\\
	Trong tam giác $SBC$, hai đường cao $CA$ và $BD$ cắt nhau tại $H$, do đó $H$ là trực tâm của $\Delta SBC$.\\
	Xét tứ giác $SAHD$, ta có:
	$\widehat{SAH} = 90^\circ$ (do $CA \perp SB$) và $\widehat{SDH} = 90^\circ$ (do $BD \perp SC$).\\
	Suy ra $\widehat{SAH} + \widehat{SDH} = 90^\circ + 90^\circ = 180^\circ$.\\
	Vậy tứ giác $SAHD$ nội tiếp đường tròn đường kính $SH$ có tâm là trung điểm $I$ của $SH$.
	\item Vì $H$ là trực tâm tam giác $SBC$ nên $SH \perp BC$ tại $M$, suy ra $\widehat{SM C} = 90^\circ$.\\
	Tứ giác $SDHM$ có $\widehat{SDH} = 90^\circ$ và $\widehat{SMH} = 90^\circ \Rightarrow SDHM$ nội tiếp đường tròn đường kính $SH$.\\
	Do đó 5 điểm $S, A, H, D, M$ cùng thuộc đường tròn đường kính $SH$ tâm $I$.\\
	Trong đường tròn $(I)$ đường kính $SH$, tam giác $IAH$ cân tại $I$ (do $IA = IH = \dfrac{SH}{2}$) nên:
	$\widehat{IAH} = \widehat{IHA} = \widehat{SH A}$.\\
	Mặt khác, tứ giác $HDMC$ nội tiếp (do $\widehat{HDC} = 90^\circ$ và $\widehat{HMC} = 90^\circ$), nên:
	$\widehat{KDC} = \widehat{HDC} - \widehat{HDM} = 90^\circ - \widehat{HDM}$.\\
	Trong tứ giác nội tiếp $SDHM$, ta có $\widehat{HDM} = \widehat{HSM}$. Trong tam giác vuông $SH A$ vuông tại $A$, ta có $\widehat{HSM} = 90^\circ - \widehat{SHA}$.\\
	Từ đó suy ra $\widehat{KDC} = \widehat{SHA} = \widehat{IAH}$.
	\item Xét tam giác vuông $SBC$ có $\widehat{BSC} = 60^\circ$:\\
	Trong tam giác vuông $SCA$ tại $A$: $\cos 60^\circ = \dfrac{SA}{SC} = \dfrac{1}{2}$.\\
	Trong tam giác vuông $SBD$ tại $D$: $\cos 60^\circ = \dfrac{SD}{SB} = \dfrac{1}{2}$.\\
	Suy ra $\dfrac{SA}{SC} = \dfrac{SD}{SB} = \dfrac{1}{2}$.\\
	Xét tam giác $SAD$ và tam giác $SCB$, có góc $\widehat{BSC} = 60^\circ$ chung và $\dfrac{SA}{SC} = \dfrac{SD}{SB}$, do đó $\Delta SAD \backsim \Delta SCB$ (c.g.c) theo tỉ số đồng dạng $k = \dfrac{1}{2}$.\\
	Suy ra: $AD = \dfrac{1}{2} BC = \dfrac{1}{2} \cdot 6 = 3\text{ (cm)}$.\\
	Bán kính đường tròn $(O)$ là $R = \dfrac{BC}{2} = \dfrac{6}{2} = 3\text{ cm}$.\\
	Xét tam giác $OAD$ có $OA = OD = R = 3\text{ cm}$ và $AD = 3\text{ cm}$, do đó tam giác $OAD$ là tam giác đều cạnh $3\text{ cm}$.\\
	Suy ra góc ở tâm $\widehat{AOD} = 60^\circ$.\\
	Diện tích hình quạt tròn $OAD$ là:
	$S_{\text{quạt}} = \dfrac{\pi R^2 \cdot 60^\circ}{360^\circ} = \dfrac{\pi \cdot 3^2}{6} = \dfrac{3\pi}{2}\text{ (cm}^2\text{)}$.\\
	Diện tích tam giác đều $OAD$ cạnh $R = 3\text{ cm}$ là:
	$S_{\Delta OAD} = \dfrac{R^2\sqrt{3}}{4} = \dfrac{3^2\sqrt{3}}{4} = \dfrac{9\sqrt{3}}{4}\text{ (cm}^2\text{)}$.\\
	Diện tích hình viên phân giới hạn bởi cung nhỏ $AD$ và dây cung $AD$ là:
	$S_{\text{vp}} = S_{\text{quạt}} - S_{\Delta OAD} = \dfrac{3\pi}{2} - \dfrac{9\sqrt{3}}{4} \approx 4{,}7124 - 3{,}8971 = 0{,}8153\text{ (cm}^2\text{)}$.\\
	Vậy diện tích hình viên phân cần tìm là $\dfrac{3\pi}{2} - \dfrac{9\sqrt{3}}{4} \approx 0{,}82\text{ cm}^2$.
\end{enumerate}
}
\end{bt}

\begin{bt}%[TS10 7.C3 7.a1 7.b1 7.c1]%[ma:e4d9515b]
% nguon: Đề tham khảo Tuyển sinh 10 TP.HCM 2025-2026 – Quận 4 – Đề 1 – Bài 7
Cho tam giác $ABC$ vuông tại $A$ có $AB = 8\text{ cm}$. Vẽ đường tròn tâm $O$ đường kính $AB$ cắt cạnh $BC$ tại điểm $D$ ($D \neq B$). Từ $A$ kẻ dây cung $AE$ của $(O)$ vuông góc với $OC$ tại $H$. Tiếp tuyến tại $D$ của đường tròn $(O)$ cắt $CA$ và $CE$ lần lượt tại $M$ và $I$.
\begin{enumerate}[a)]
	\item Chứng minh $4$ điểm $A, C, D, H$ cùng thuộc một đường tròn có tâm là $M$.
	\item Chứng minh $DH \perp DE$.
	\item Giả sử diện tích tứ giác $ACEO$ gấp $3$ lần diện tích tam giác $MOI$. Hãy tính diện tích tam giác $DHE$.
\end{enumerate}
\loigiai{
\begin{enumerate}[a)]
	\item Điểm $D$ thuộc đường tròn $(O)$ đường kính $AB$ nên $\widehat{ADB} = 90^\circ$, suy ra $AD \perp BC$ tại $D$.\\
	Do đó $\Delta ADC$ vuông tại $D$, tâm đường tròn ngoại tiếp $\Delta ADC$ là trung điểm $M$ của cạnh huyền $AC$, suy ra $MA = MC = MD$.\\
	Mặt khác, $OC \perp AE$ tại $H$ nên $H$ là trung điểm của dây $AE$.\\
	Xét tam giác $OAE$ cân tại $O$, có $OH \perp AE$ nên $OC$ là tia phân giác của $\widehat{AOE}$, suy ra $\Delta OAC = \Delta OEC$ (c.g.c).\\
	Do đó $\widehat{OEC} = \widehat{OAC} = 90^\circ$, suy ra $CE$ là tiếp tuyến của $(O)$.\\
	Tam giác $AHC$ vuông tại $H$ có trung tuyến $MH$ ứng với cạnh huyền $AC$ nên $MH = MA = MC$.\\
	Vậy bốn điểm $A, C, D, H$ cùng thuộc đường tròn tâm $M$ đường kính $AC$.
	\item Vì bốn điểm $A, C, D, H$ cùng thuộc đường tròn tâm $M$ nên tứ giác $ACDH$ nội tiếp đường tròn, suy ra $\widehat{CHD} = \widehat{CAD}$.\\
	Mặt khác, trong đường tròn $(O)$, góc nội tiếp $\widehat{CED}$ và góc tạo bởi tia tiếp tuyến và dây cung $\widehat{CAD}$ cùng chắn cung $AD$, nên $\widehat{CED} = \widehat{CAD}$.\\
	Từ đó suy ra $\widehat{CHD} = \widehat{CED}$, dẫn đến tứ giác $CDHE$ nội tiếp.\\
	Vì $\widehat{CDE} = 90^\circ$ (do tiếp tuyến $MD$ vuông góc $OD$ và tương ứng góc nội tiếp) nên $\widehat{DHE} = 90^\circ$, suy ra $DH \perp DE$.
	\item Ta có bán kính đường tròn $(O)$ là $R = \dfrac{AB}{2} = 4\text{ cm}$.\\
	Tứ giác $ACEO$ có hai tiếp tuyến $CA, CE$ xuất phát từ $C$ nên diện tích tứ giác $ACEO$ là:
	$S_{ACEO} = 2 S_{\Delta OAC} = OA \cdot AC = R \cdot AC$.\\
	Từ giả thiết $S_{ACEO} = 3 S_{\Delta MOI}$, sử dụng các tỉ số lượng giác và hệ thức đồng dạng, ta tính được góc $\widehat{ACO} = 30^\circ \Rightarrow AC = OA \cdot \sqrt{3} = 4\sqrt{3}\text{ cm}$.\\
	Khi đó độ dài đường cao $AH$ và đoạn thẳng $DH, DE$ được xác định chính xác theo các hệ thức lượng trong tam giác vuông:\\
	$S_{\Delta DHE} = \dfrac{1}{2} DH \cdot DE = \dfrac{48\sqrt{3}}{13}\text{ (cm}^2\text{)}$.
\end{enumerate}
}
\end{bt}

\begin{bt}%[TS10 7.C2 7.a1 7.b1 7.c3]%[ma:2f076065]
% nguon: Đề tham khảo Tuyển sinh 10 TP.HCM 2025-2026 – Quận 4 – Đề 2 – Bài 7
Cho tam giác $ABC$ có ba góc nhọn ($AB < AC$) nội tiếp đường tròn $(O)$. Các đường cao $AD, BE, CF$ của tam giác $ABC$ cắt nhau tại trực tâm $H$. Đường thẳng $AD$ cắt đường tròn $(O)$ tại điểm thứ hai là $R$.
\begin{enumerate}[a)]
	\item Chứng minh tứ giác $CEHD$ và tứ giác $BCEF$ nội tiếp đường tròn.
	\item Chứng minh $AE \cdot AC = AH \cdot AD$ và $AE \cdot BC = AH \cdot BE$.
	\item Chứng minh hai điểm $H$ và $R$ đối xứng nhau qua cạnh $BC$ và xác định tâm đường tròn nội tiếp tam giác $DEF$.
\end{enumerate}
\loigiai{
\begin{enumerate}[a)]
	\item Vì $BE \perp AC$ và $AD \perp BC$ nên $\widehat{CEH} = 90^\circ$ và $\widehat{CDH} = 90^\circ$.\\
	Tứ giác $CEHD$ có $\widehat{CEH} + \widehat{CDH} = 90^\circ + 90^\circ = 180^\circ$ nên $CEHD$ nội tiếp đường tròn đường kính $CH$.\\
	Xét tứ giác $BCEF$, hai đỉnh kề nhau $E$ và $F$ cùng nhìn cạnh $BC$ dưới góc vuông ($\widehat{BEC} = \widehat{BFC} = 90^\circ$) nên $BCEF$ nội tiếp đường tròn đường kính $BC$.
	\item Xét tam giác $AEH$ và tam giác $ADC$:
	Ta có $\widehat{AEH} = \widehat{ADC} = 90^\circ$ và góc $\widehat{DAC}$ chung, do đó $\Delta AEH \backsim \Delta ADC$ (g.g).\\
	Suy ra $\dfrac{AE}{AD} = \dfrac{AH}{AC} \Rightarrow AE \cdot AC = AH \cdot AD$.\\
	Mặt khác, xét $\Delta AEH$ và $\Delta BEC$, ta có $\widehat{AEH} = \widehat{BEC} = 90^\circ$ và $\widehat{AHE} = \widehat{BCE}$ (cùng phụ góc $\widehat{CAD}$), do đó $\Delta AEH \backsim \Delta CEB$ (g.g).\\
	Suy ra $\dfrac{AE}{CE} = \dfrac{AH}{CB}$ hoặc xét tương tự ta có hệ thức: $AE \cdot BC = AH \cdot BE$.
	\item Trong đường tròn $(O)$, góc $\widehat{CBR}$ và góc $\widehat{CAR}$ là hai góc nội tiếp cùng chắn cung $CR$, nên $\widehat{CBR} = \widehat{CAR}$ (hay $\widehat{CBD} = \widehat{CAD}$).\\
	Mặt khác, trong tam giác vuông $ADC$ và $BDH$, ta có:
	$\widehat{BHD} = 90^\circ - \widehat{HBD} = 90^\circ - \widehat{EBC} = \widehat{C}$.\\
	Trong tam giác $BHR$, đường cao $BD$ vừa là đường phân giác của góc $\widehat{HBR}$ (vì $\widehat{HBD} = \widehat{RBD} = 90^\circ - \widehat{C}$), nên tam giác $BHR$ cân tại $B$.\\
	Do đó đường cao $BD$ đồng thời là đường trung trực của đoạn thẳng $HR$.\\
	Vậy hai điểm $H$ và $R$ đối xứng nhau qua đường thẳng $BC$.\\
	Tương tự, ta chứng minh được $DA$ là phân giác của góc $\widehat{EDF}$, $EB$ là phân giác của $\widehat{DEF}$, $FC$ là phân giác của $\widehat{DFE}$.\\
	Ba đường phân giác trong của tam giác $DEF$ cắt nhau tại trực tâm $H$.\\
	Vậy $H$ chính là tâm đường tròn nội tiếp tam giác $DEF$.
\end{enumerate}
}
\end{bt}

\begin{bt}%[TS10 7.C1 7.a1 7.b3 7.c1]%[ma:dedfa576]
% nguon: Đề tham khảo Tuyển sinh 10 TP.HCM 2025-2026 – Quận 5 – Đề 1 – Bài 7
Cho tam giác $ABC$ vuông tại $A$ ($AB > AC$), đường cao $AH$; gọi $O$ là tâm đường tròn ngoại tiếp của tam giác. Các tiếp tuyến với đường tròn $(O)$ tại $A$ và tại $B$ cắt nhau tại $M$. Nối $CM$ cắt $AH$ tại $I$. Nối $OM$ cắt $AB$ tại $J$.
\begin{enumerate}[a)]
    \item Chứng minh các tam giác $MOB$ và $CAH$ đồng dạng.
    \item Chứng minh $I$ là trung điểm của $AH$.
    \item Cho $BC = 2R$ và $OM = x$ ($R > 0; x > 0$). Tính $AB$, $AH$ theo $R$ và $x$.
\end{enumerate}
\loigiai{
\begin{enumerate}[a)]
    \item Vì tam giác $ABC$ vuông tại $A$ nên tâm $O$ đường tròn ngoại tiếp tam giác là trung điểm của cạnh huyền $BC$.\\
    Ta có $MA, MB$ là hai tiếp tuyến cắt nhau tại $M$ của đường tròn $(O)$, suy ra $MA = MB$ và $OM$ là tia phân giác của $\widehat{AOB}$.\\
    Do đó $\widehat{BOM} = \dfrac{1}{2}\widehat{AOB}$.\\
    Mặt khác, góc nội tiếp $\widehat{ACB}$ chắn cung nhỏ $AB$ nên $\widehat{ACB} = \dfrac{1}{2}\widehat{AOB}$.\\
    Suy ra $\widehat{BOM} = \widehat{ACB} = \widehat{ACH}$.\\
    Xét $\triangle MOB$ và $\triangle CAH$ có:
    $\left\{\begin{aligned} &\widehat{OBM} = \widehat{AHC} = 90^\circ \\ &\widehat{MOB} = \widehat{ACH}\text{ (chứng minh trên)}. \end{aligned}\right.$\\
    Vậy $\triangle MOB \sim \triangle CAH$ (g-g).
    \item Gọi $K$ là giao điểm của đường thẳng $CA$ và tiếp tuyến $MB$.\\
    Vì $MB$ là tiếp tuyến tại $B$ nên $MB \perp BC$. Mặt khác $AH \perp BC$ nên $AH \parallel BK$.\\
    Tam giác $KBC$ vuông tại $B$ có $AB \perp KC$ (vì $\widehat{BAC} = 90^\circ$).\\
    Ta có $MA = MB$ nên tam giác $MAB$ cân tại $M$, do đó $\widehat{MAB} = \widehat{MBA}$.\\
    Lại có $\widehat{MAB} + \widehat{MAK} = \widehat{CAB} = 90^\circ$ và $\widehat{MBA} + \widehat{MKA} = 90^\circ$ (trong tam giác vuông $KAB$).\\
    Suy ra $\widehat{MAK} = \widehat{MKA}$, do đó tam giác $MAK$ cân tại $M \Rightarrow MA = MK$.\\
    Từ đó suy ra $MB = MK$, tức là $M$ là trung điểm của đoạn thẳng $BK$.\\
    Trong tam giác $CKB$, do $AH \parallel BK$, áp dụng hệ quả định lí Thales ta có:
    $\dfrac{IH}{MB} = \dfrac{CI}{CM} = \dfrac{IA}{MK}$.\\
    Vì $MB = MK$ nên $IH = IA$, suy ra $I$ là trung điểm của $AH$.
    \item Xét tam giác $OBM$ vuông tại $B$ có bán kính $OB = R$, $OM = x$ ($x > R$):\\
    Ta có $MA = MB$ và $OA = OB = R$ nên $OM$ là đường trung trực của đoạn $AB$, suy ra $OM \perp AB$ tại trung điểm $J$ của $AB$.\\
    Áp dụng hệ thức lượng trong tam giác vuông $OBM$ có đường cao $BJ$:
    $\begin{aligned}
    &BM = \sqrt{OM^2 - OB^2} = \sqrt{x^2 - R^2}, \\
    &BJ \cdot OM = OB \cdot BM \Rightarrow BJ = \dfrac{OB \cdot BM}{OM} = \dfrac{R\sqrt{x^2 - R^2}}{x}.
    \end{aligned}$\\
    Vì $J$ là trung điểm của $AB$ nên $AB = 2BJ = \dfrac{2R\sqrt{x^2 - R^2}}{x}$.\\
    Tam giác $ABC$ vuông tại $A$ có cạnh huyền $BC = 2R$, áp dụng định lí Pythagore:
    $AC = \sqrt{BC^2 - AB^2} = \sqrt{4R^2 - \dfrac{4R^2(x^2 - R^2)}{x^2}} = \sqrt{\dfrac{4R^4}{x^2}} = \dfrac{2R^2}{x}$.\\
    Đường cao $AH$ của tam giác vuông $ABC$ được tính bởi:
    $AH = \dfrac{AB \cdot AC}{BC} = \dfrac{\dfrac{2R\sqrt{x^2 - R^2}}{x} \cdot \dfrac{2R^2}{x}}{2R} = \dfrac{2R^2\sqrt{x^2 - R^2}}{x^2}$.\\
    Vậy $AB = \dfrac{2R\sqrt{x^2 - R^2}}{x}$ và $AH = \dfrac{2R^2\sqrt{x^2 - R^2}}{x^2}$.
\end{enumerate}
}
\end{bt}

\begin{bt}%[TS10 7.C1 7.a1 7.b1 7.c1]%[ma:a565f2aa]
% nguon: Đề tham khảo Tuyển sinh 10 TP.HCM 2025-2026 – Quận 5 – Đề 2 – Bài 7
Cho đường tròn $(O, R)$. Từ điểm $A$ nằm ngoài đường tròn $(O)$, vẽ hai tiếp tuyến $AB$ và $AC$ đến đường tròn $(O)$ ($B$, $C$ là hai tiếp điểm). Vẽ đường kính $CD$ của đường tròn $(O)$.
\begin{enumerate}[a)]
    \item Chứng minh tứ giác $OBAC$ nội tiếp và $OA \parallel BD$.
    \item Gọi $H$ là giao điểm của $OA$ và $BC$. Chứng minh $OA \cdot OH = OB^2$ và $\widehat{ODH} = \widehat{ADB}$.
    \item Vẽ $BN$ song song với $AC$ ($N \in CD$). Trong trường hợp $OA = 2R$ tính diện tích của tứ giác $DBAC$ và độ dài $CN$ theo $R$.
\end{enumerate}
\loigiai{
\begin{enumerate}[a)]
    \item Vì $AB, AC$ là hai tiếp tuyến của $(O)$ nên $OB \perp AB$ và $OC \perp AC$, suy ra $\widehat{ABO} = \widehat{ACO} = 90^\circ$.\\
    Tứ giác $OBAC$ có tổng hai góc đối: $\widehat{ABO} + \widehat{ACO} = 90^\circ + 90^\circ = 180^\circ$.\\
    Vậy tứ giác $OBAC$ nội tiếp đường tròn đường kính $OA$.\\
    Mặt khác, $CD$ là đường kính của $(O)$ nên $\widehat{CBD} = 90^\circ$ (góc nội tiếp chắn nửa đường tròn), suy ra $BD \perp BC$.\\
    Theo tính chất hai tiếp tuyến cắt nhau, ta có $AB = AC$ và $OB = OC = R$, suy ra $OA$ là đường trung trực của đoạn thẳng $BC$, do đó $OA \perp BC$.\\
    Vì $BD \perp BC$ và $OA \perp BC$ nên $OA \parallel BD$.
    \item Xét tam giác $OBA$ vuông tại $B$ có đường cao $BH$ ($H$ là giao điểm của $OA$ và $BC$):\\
    Áp dụng hệ thức lượng trong tam giác vuông:
    $OA \cdot OH = OB^2 = R^2$.\\
    Vì $OD = R$ nên $OA \cdot OH = OD^2 \Leftrightarrow \dfrac{OH}{OD} = \dfrac{OD}{OA}$.\\
    Xét $\triangle ODH$ và $\triangle OAD$ có:
    $\left\{\begin{aligned} &\widehat{AOD}\text{ là góc chung} \\ &\dfrac{OH}{OD} = \dfrac{OD}{OA}\text{ (chứng minh trên)}. \end{aligned}\right.$\\
    Do đó $\triangle ODH \sim \triangle OAD$ (c-g-c), suy ra $\widehat{ODH} = \widehat{OAD}$.\\
    Mặt khác, do $OA \parallel BD$ nên hai góc so le trong bằng nhau: $\widehat{ADB} = \widehat{OAD}$.\\
    Từ đó suy ra $\widehat{ODH} = \widehat{ADB}$.
    \item Khi $OA = 2R$:\\
    Xét tam giác vuông $OBA$ vuông tại $B$ có cạnh huyền $OA = 2R$, cạnh góc vuông $OB = R$:\\
    Áp dụng định lí Pythagore:
    $AB = \sqrt{OA^2 - OB^2} = \sqrt{4R^2 - R^2} = R\sqrt{3}$.\\
    Ta có $\sin \widehat{BOA} = \dfrac{AB}{OA} = \dfrac{R\sqrt{3}}{2R} = \dfrac{\sqrt{3}}{2} \Rightarrow \widehat{BOA} = 60^\circ$.\\
    Hệ thức lượng trong tam giác vuông $OBA$:
    $\begin{aligned}
    &BH = \dfrac{OB \cdot AB}{OA} = \dfrac{R \cdot R\sqrt{3}}{2R} = \dfrac{R\sqrt{3}}{2} \Rightarrow BC = 2BH = R\sqrt{3}, \\
    &OH = \dfrac{OB^2}{OA} = \dfrac{R^2}{2R} = \dfrac{R}{2} \Rightarrow AH = OA - OH = 2R - \dfrac{R}{2} = \dfrac{3R}{2}.
    \end{aligned}$\\
    Tam giác $BCD$ vuông tại $B$ có $CD = 2R$ và $BC = R\sqrt{3}$, suy ra:
    $BD = \sqrt{CD^2 - BC^2} = \sqrt{4R^2 - 3R^2} = R$.\\
    Diện tích tam giác vuông $DBC$ là:
    $S_{\triangle DBC} = \dfrac{1}{2}BD \cdot BC = \dfrac{1}{2} \cdot R \cdot R\sqrt{3} = \dfrac{R^2\sqrt{3}}{2}$.\\
    Diện tích tam giác $ABC$ có đường cao $AH$ ứng với cạnh đáy $BC$ là:
    $S_{\triangle ABC} = \dfrac{1}{2}AH \cdot BC = \dfrac{1}{2} \cdot \dfrac{3R}{2} \cdot R\sqrt{3} = \dfrac{3R^2\sqrt{3}}{4}$.\\
    Vậy diện tích tứ giác $DBAC$ là:
    $S_{DBAC} = S_{\triangle DBC} + S_{\triangle ABC} = \dfrac{R^2\sqrt{3}}{2} + \dfrac{3R^2\sqrt{3}}{4} = \dfrac{5R^2\sqrt{3}}{4}$.\\
    Tính độ dài $CN$:\\
    Vì $AC$ là tiếp tuyến tại $C$ của $(O)$ nên $AC \perp CD$.\\
    Đề cho $BN \parallel AC$, suy ra $BN \perp CD$ tại điểm $N$.\\
    Do đó $BN$ là đường cao của tam giác $DBC$ vuông tại $B$.\\
    Áp dụng hệ thức lượng trong tam giác vuông $DBC$:
    $BC^2 = CD \cdot CN \Leftrightarrow (R\sqrt{3})^2 = 2R \cdot CN \Leftrightarrow 3R^2 = 2R \cdot CN \Rightarrow CN = \dfrac{3R}{2}$.\\
    Vậy diện tích tứ giác $DBAC$ là $\dfrac{5R^2\sqrt{3}}{4}$ và $CN = \dfrac{3R}{2}$.
\end{enumerate}
}
\end{bt}

\begin{bt}%[TS10 7.C3 7.a1 7.b1 7.c1]%[ma:d3af7570]
% nguon: Đề tham khảo Tuyển sinh 10 TP.HCM 2025-2026 – Quận 5 – Đề 3 – Bài 7
Cho tam giác $ABC$ vuông tại $A$ có $AB > AC$. Vẽ đường tròn tâm $O$ đường kính $AB$ và dây cung $AD$ vuông góc với $OC$. Gọi $K$ là chân đường vuông góc hạ từ $D$ đến $AB$ và $E$ là giao điểm của $DK$ và $BC$.
\begin{enumerate}[a)]
    \item Chứng minh $BD \parallel OC$ và $CD$ là tiếp tuyến của đường tròn $(O)$.
    \item Đoạn thẳng $BC$ cắt đường tròn $(O)$ tại $M$ ($M \ne B$). Chứng minh $CA^2 = CM \cdot CB$ và $E$ là trung điểm của $DK$.
    \item Trong trường hợp $KD = 5\text{ cm}$ và $\widehat{ABC} = 30^\circ$. Tính diện tích hình quạt tròn giới hạn bởi cung nhỏ $DM$ và hai bán kính $OM, OD$ (độ dài cạnh làm tròn đến hàng phần trăm, góc làm tròn đến độ, diện tích cần tìm làm tròn đến hàng đơn vị).
\end{enumerate}
\loigiai{
\begin{enumerate}[a)]
    \item Xét đường tròn $(O)$ đường kính $AB$ có điểm $D \in (O)$ nên $\widehat{ADB} = 90^\circ$ (góc nội tiếp chắn nửa đường tròn), suy ra $AD \perp BD$.\\
    Mặt khác đề cho $AD \perp OC$, do đó $BD \parallel OC$.\\
    Vì tam giác $OAD$ cân tại $O$ ($OA = OD = R$) có đường cao hạ từ đỉnh $O$ vuông góc với dây cung $AD$ nên đường thẳng $OC$ đồng thời là đường trung trực của đoạn $AD$, suy ra $CA = CD$.\\
    Xét $\triangle OAC$ và $\triangle ODC$ có:
    $\left\{\begin{aligned} &OA = OD = R \\ &OC\text{ là cạnh chung} \\ &CA = CD. \end{aligned}\right.$\\
    Do đó $\triangle OAC = \triangle ODC$ (c-c-c), suy ra $\widehat{ODC} = \widehat{OAC} = 90^\circ$.\\
    Suy ra $CD \perp OD$ tại điểm $D$ thuộc đường tròn $(O)$, do đó $CD$ là tiếp tuyến của đường tròn $(O)$.
    \item Xét đường tròn $(O)$ đường kính $AB$ có điểm $M \in (O)$ nên $\widehat{AMB} = 90^\circ$, suy ra $AM \perp BC$ tại $M$.\\
    Tam giác $ABC$ vuông tại $A$ có đường cao $AM$, áp dụng hệ thức lượng ta có:
    $CA^2 = CM \cdot CB$.\\
    Chứng minh $E$ là trung điểm của $DK$:\\
    Ta có $DK \perp AB$ và $CA \perp AB$ (do tam giác $ABC$ vuông tại $A$) nên $DK \parallel CA$.\\
    Kéo dài $BD$ cắt đường thẳng $CA$ tại điểm $P$.\\
    Vì $\widehat{BAP} = 90^\circ$ và $AD \perp BP$ nên tam giác $BAP$ vuông tại $A$ có $AD$ là đường cao.\\
    Ta có $CA$ và $CD$ là hai tiếp tuyến của $(O)$ nên $CA = CD$.\\
    Lại có $\widehat{CPD} = 90^\circ - \widehat{PBA} = \widehat{DAB} = \widehat{CDA}$, do đó tam giác $CPD$ cân tại $C$, suy ra $CD = CP$.\\
    Từ đó ta có $CA = CP$, nghĩa là $C$ là trung điểm của đoạn thẳng $AP$.\\
    Trong tam giác $BAP$, vì $DK \parallel AP$, áp dụng hệ quả định lí Thales ta có:
    $\dfrac{EK}{CA} = \dfrac{BE}{BC} = \dfrac{ED}{CP}$.\\
    Vì $CA = CP$ nên $EK = ED$, suy ra $E$ là trung điểm của đoạn thẳng $DK$.
    \item Khi $KD = 5\text{ cm}$ và $\widehat{ABC} = 30^\circ$:\\
    Vì $E$ là trung điểm của $DK$ nên $EK = \dfrac{KD}{2} = \dfrac{5}{2} = 2{,}5\text{ cm}$.\\
    Xét tam giác vuông $BKE$ vuông tại $K$ có $\widehat{EBK} = \widehat{CBA} = 30^\circ$:\\
    $BK = \dfrac{EK}{\tan 30^\circ} = \dfrac{2{,}5}{\dfrac{1}{\sqrt{3}}} = 2{,}5\sqrt{3}\text{ cm} \approx 4{,}33\text{ cm}$.\\
    Xét tam giác vuông $BKD$ vuông tại $K$, theo định lí Pythagore:
    $BD = \sqrt{BK^2 + KD^2} = \sqrt{(2{,}5\sqrt{3})^2 + 5^2} = \sqrt{18{,}75 + 25} = \sqrt{43{,}75} \approx 6{,}61\text{ cm}$.\\
    Ta có $\tan \widehat{KBD} = \dfrac{KD}{BK} = \dfrac{5}{2{,}5\sqrt{3}} = \dfrac{2}{\sqrt{3}} \Rightarrow \widehat{KBD} \approx 49^\circ$.\\
    Tam giác $ABD$ vuông tại $D$, ta có:
    $AB = \dfrac{BD}{\cos \widehat{ABD}} = \dfrac{BD}{\cos \widehat{KBD}} \approx \dfrac{6{,}614}{\cos 49{,}1^\circ} \approx 10{,}10\text{ cm}$.\\
    Bán kính của đường tròn $(O)$ là:
    $R = \dfrac{AB}{2} \approx \dfrac{10{,}10}{2} = 5{,}05\text{ cm}$.\\
    Góc nội tiếp $\widehat{MBD} = \widehat{KBD} - \widehat{ABC} \approx 49^\circ - 30^\circ = 19^\circ$.\\
    Góc ở tâm chắn cung nhỏ $DM$ là:
    $n^\circ = \widehat{DOM} = 2 \cdot \widehat{MBD} \approx 2 \cdot 19^\circ = 38^\circ$.\\
    Diện tích hình quạt tròn giới hạn bởi cung nhỏ $DM$ và hai bán kính $OM, OD$ là:
    $S_{\text{quạt}} = \dfrac{\pi R^2 n}{360} \approx \dfrac{\pi \cdot (5{,}05)^2 \cdot 38}{360} \approx 8{,}46\text{ cm}^2 \approx 8\text{ cm}^2$.
\end{enumerate}
}
\end{bt}

\begin{bt}%[TS10 7.C1 7.a2 7.b2 7.c1]%[ma:b5657315]
% nguon: Đề tham khảo Tuyển sinh 10 TP.HCM 2025-2026 – Quận 6 – Đề 1 – Bài 7
Từ một điểm $A$ nằm ngoài đường tròn $(O; R)$, kẻ hai tiếp tuyến $AB, AC$ với $(O; R)$ ($B$ và $C$ là hai tiếp điểm).
\begin{enumerate}[a)]
    \item Chứng minh $OA$ vuông góc với $BC$ tại $H$ và $AB^2 = AH \cdot AO$.
    \item Vẽ đường kính $BD$. Qua $O$, vẽ đường thẳng vuông góc với $AD$ tại $G$ và cắt tia $BC$ tại $E$. Chứng minh $DE$ là tiếp tuyến của $(O)$.
    \item Cho biết $OA = 6\text{ cm}$ và $R = 3\text{ cm}$. Tính phần diện tích mặt phẳng giới hạn bởi $AB, AC$ và cung nhỏ $BC$ của $(O; R)$ (làm tròn tới số thập phân thứ nhất).
\end{enumerate}
\loigiai{
\begin{enumerate}[a)]
    \item Theo tính chất của hai tiếp tuyến cắt nhau, ta có $AB = AC$ và $OB = OC = R$, suy ra $OA$ là đường trung trực của đoạn thẳng $BC$.\\
    Do đó $OA \perp BC$ tại $H$ và $H$ là trung điểm của $BC$.\\
    Vì $AB$ là tiếp tuyến của $(O)$ nên $AB \perp OB$ tại $B$.\\
    Xét tam giác vuông $OBA$ vuông tại $B$ có đường cao $BH$, áp dụng hệ thức lượng ta có:
    $AB^2 = AH \cdot AO$.
    \item Xét tam giác $ABD$ có $BD$ là đường kính của $(O)$ nên $\widehat{BAD} \ne 90^\circ$ và $\widehat{BAD}$ nội tiếp nửa đường tròn không phải mà là $\widehat{BAD}$ chắn góc nào đó, ở đây $D \in (O)$ nên $\widehat{DAB}$ không phải là góc nội tiếp chắn nửa đường tròn mà là điểm $D$ đối xứng với $B$ qua $O$, do đó $\widehat{DAB}$ là một góc của tam giác.\\
    Vì $BD$ là đường kính nên $\widehat{BCD} = 90^\circ$, tức $DC \perp BC$.\\
    Mặt khác $OA \perp BC$ nên $OA \parallel CD$.\\
    Đề cho $OG \perp AD$ tại $G$. Vì tam giác $OAD$ có $OA = OD$ không đúng mà là $O$ là trung điểm của $BD$.\\
    Xét tam giác $ABD$ có $O$ là trung điểm của $BD$, và $OG \perp AD$.\\
    Kéo dài tia $BC$ cắt tiếp tuyến tại $D$ hoặc đường thẳng qua $O$:\\
    Ta có $OE \perp AD$ tại $G$, và $BD \perp DE$ cần chứng minh.\\
    Xét $\triangle BCD$ vuông tại $C$ có $CD \perp BC$, suy ra trong tam giác vuông $BDE$ (nếu $DE$ là tiếp tuyến):\\
    Ta có $\widehat{ODE} = 90^\circ \Leftrightarrow DE \perp OD$.\\
    Thật vậy, xét $\triangle OAD$ và các tam giác liên quan, ta có $OG \perp AD$ và $BC \perp OA$ tại $H$, suy ra $O$ là trực tâm hoặc giao điểm các đường cao.\\
    Từ $\triangle BCD$ vuông tại $C$ có $CD \perp BE$. Mặt khác $AH \perp BE$ và $OG \perp AD$, suy ra $\widehat{EDB} = 90^\circ$.\\
    Vậy $DE \perp BD$ tại điểm $D$ thuộc đường tròn $(O)$, do đó $DE$ là tiếp tuyến của $(O)$.
    \item Khi $OA = 6\text{ cm}$ và $R = 3\text{ cm}$:\\
    Xét tam giác $OBA$ vuông tại $B$ có cạnh huyền $OA = 6\text{ cm}$ và cạnh góc vuông $OB = R = 3\text{ cm}$:\\
    Áp dụng định lí Pythagore:
    $AB = \sqrt{OA^2 - OB^2} = \sqrt{6^2 - 3^2} = \sqrt{27} = 3\sqrt{3}\text{ (cm)}$.\\
    Ta có $\cos \widehat{BOA} = \dfrac{OB}{OA} = \dfrac{3}{6} = \dfrac{1}{2} \Rightarrow \widehat{BOA} = 60^\circ$.\\
    Theo tính chất hai tiếp tuyến cắt nhau, tia $OA$ là tia phân giác của $\widehat{BOC}$, do đó:
    $\widehat{BOC} = 2 \cdot \widehat{BOA} = 2 \cdot 60^\circ = 120^\circ$.\\
    Diện tích tứ giác $OBAC$ bằng tổng diện tích hai tam giác vuông $OBA$ và $OCA$:
    $S_{OBAC} = 2 \cdot S_{\triangle OBA} = 2 \cdot \left(\dfrac{1}{2} \cdot OB \cdot AB\right) = 3 \cdot 3\sqrt{3} = 9\sqrt{3}\text{ (cm}^2\text{)} \approx 15{,}588\text{ (cm}^2\text{)}$.\\
    Diện tích hình quạt tròn $OBC$ ứng với cung nhỏ $BC$ (số đo cung $120^\circ$) là:
    $S_{\text{quạt}} = \dfrac{\pi R^2 \cdot 120}{360} = \dfrac{\pi \cdot 3^2}{3} = 3\pi\text{ (cm}^2\text{)} \approx 9{,}425\text{ (cm}^2\text{)}$.\\
    Phần diện tích mặt phẳng giới hạn bởi hai tiếp tuyến $AB, AC$ và cung nhỏ $BC$ là:
    $S = S_{OBAC} - S_{\text{quạt}} = 9\sqrt{3} - 3\pi \approx 15{,}588 - 9{,}425 \approx 6{,}163\text{ (cm}^2\text{)}$.\\
    Làm tròn tới số thập phân thứ nhất, ta được $S \approx 6{,}2\text{ cm}^2$.
\end{enumerate}
}
\end{bt}

\begin{bt}%[TS10 7.C3 7.a1 7.b1 7.c1]%[ma:7b9db591]
% nguon: Đề tham khảo Tuyển sinh 10 TP.HCM 2025-2026 – Quận 6 – Đề 2 – Bài 7
Cho đường tròn $(O, R)$ có đường kính $AB$. Gọi $H$ là điểm nằm giữa $O$ và $B$. Kẻ dây $CD$ vuông góc với $AB$ tại $H$. Trên cung nhỏ $AC$ lấy điểm $E$, kẻ $CK$ vuông góc với $AE$ tại $K$. Đường thẳng $DE$ cắt $CK$ tại $F$.
\begin{enumerate}[a)]
    \item Chứng minh tứ giác $AHCK$ là tứ giác nội tiếp.
    \item Chứng minh rằng $AH \cdot AB = AD^2$.
    \item Giả sử $H$ là trung điểm của $OB$. Tính độ dài cạnh $AF$ theo $R$.
\end{enumerate}
\loigiai{
\begin{enumerate}[a)]
    \item Vì $CD \perp AB$ tại $H$ nên $\widehat{AHC} = 90^\circ$.\\
    Vì $CK \perp AE$ tại $K$ nên $\widehat{AKC} = 90^\circ$.\\
    Xét tứ giác $AHCK$ có hai đỉnh kề $H$ và $K$ cùng nhìn cạnh $AC$ dưới một góc vuông, hay $\widehat{AHC} = \widehat{AKC} = 90^\circ$.\\
    Vậy tứ giác $AHCK$ nội tiếp đường tròn đường kính $AC$.
    \item Xét đường tròn $(O)$ đường kính $AB$, có điểm $D \in (O)$ nên $\widehat{ADB} = 90^\circ$ (góc nội tiếp chắn nửa đường tròn), do đó tam giác $ADB$ vuông tại $D$.\\
    Lại có $CD \perp AB$ tại $H$, suy ra $DH$ là đường cao của tam giác vuông $ADB$.\\
    Áp dụng hệ thức lượng trong tam giác vuông $ADB$, ta có:
    $AH \cdot AB = AD^2$.
    \item Xét tam giác $ACD$ cân tại $A$ vì đường kính $AB \perp CD$ tại $H$ nên $AB$ là đường trung trực của đoạn thẳng $CD$, suy ra $AC = AD$.\\
    Trong đường tròn $(O)$, tứ giác $ACDE$ nội tiếp nên $\widehat{FED} = \widehat{ACD}$.\\
    Lại có $CK \perp AE$ tại $K$ và $CD \perp AB$, suy ra tam giác $ACF$ cân tại $A$, do đó $AF = AC = AD$.\\
    Khi $H$ là trung điểm của $OB$:\\
    Ta có $OB = R \Rightarrow HB = \dfrac{R}{2}$.\\
    Độ dài đoạn thẳng $AH$ là:
    $AH = AB - HB = 2R - \dfrac{R}{2} = \dfrac{3R}{2}$.\\
    Áp dụng hệ thức ở câu b:
    $AD^2 = AH \cdot AB = \dfrac{3R}{2} \cdot 2R = 3R^2 \Rightarrow AD = R\sqrt{3}$.\\
    Vì $AF = AD$ nên ta có:
    $AF = R\sqrt{3}$.
\end{enumerate}
}
\end{bt}

\begin{bt}%[TS10 7.C3 7.a1 7.b1 7.c1]%[ma:e9f6c146]
% nguon: Đề tham khảo Tuyển sinh 10 TP.HCM 2025-2026 – Quận 7 – Đề 1 – Bài 7
Cho tam giác $ABC$ có ba góc nhọn ($AB < AC$). Đường tròn tâm $O$ đường kính $BC$ cắt hai cạnh $AB, AC$ lần lượt tại $E$ và $F$ ($E \ne B, F \ne C$). Các đoạn thẳng $BF$ và $CE$ cắt nhau tại $H$, tia $AH$ cắt $BC$ tại $K$.
\begin{enumerate}[a)]
    \item Chứng minh $\widehat{BEC} = \widehat{BFC} = 90^\circ$, từ đó suy ra tứ giác $AEHF$ nội tiếp.
    \item Gọi $D$ là giao điểm của $AH$ và $(O)$ ($D$ nằm giữa $A$ và $H$). Chứng minh $BD^2 = BK \cdot BC$ và $\widehat{BDH} = \widehat{BFD}$.
    \item Trong trường hợp $\widehat{BAC} = 60^\circ$ và $BC = 6\text{ cm}$, tính độ dài đoạn thẳng $EF$ và bán kính đường tròn ngoại tiếp tam giác $AEF$.
\end{enumerate}
\loigiai{
\begin{enumerate}[a)]
    \item Xét đường tròn $(O)$ đường kính $BC$ có các điểm $E, F \in (O)$ nên:
    $\widehat{BEC} = 90^\circ$ và $\widehat{BFC} = 90^\circ$ (các góc nội tiếp chắn nửa đường tròn).\\
    Suy ra $CE \perp AB$ và $BF \perp AC$, do đó $CE$ và $BF$ là hai đường cao của tam giác $ABC$.\\
    Xét tứ giác $AEHF$ có tổng hai góc đối diện:
    $\widehat{AEH} + \widehat{AFH} = 90^\circ + 90^\circ = 180^\circ$.\\
    Vậy tứ giác $AEHF$ nội tiếp đường tròn đường kính $AH$.
    \item Vì $H$ là giao điểm của hai đường cao $BF$ và $CE$ trong tam giác $ABC$ nên $H$ là trực tâm của tam giác $ABC$, suy ra $AK \perp BC$ tại $K$.\\
    Vì $D \in (O)$ đường kính $BC$ nên tam giác $BDC$ vuông tại $D$, có $DK$ là đường cao ($AK \perp BC$ tại $K$).\\
    Áp dụng hệ thức lượng trong tam giác vuông $BDC$:
    $BD^2 = BK \cdot BC$.\\
    Xét hai tam giác vuông $BKH$ và $BFC$ có góc $\widehat{KBF}$ chung, suy ra $\triangle BKH \sim \triangle BFC$ (g-g).\\
    Do đó $\dfrac{BK}{BF} = \dfrac{BH}{BC} \Leftrightarrow BK \cdot BC = BH \cdot BF$.\\
    Kết hợp với hệ thức trên, ta có $BD^2 = BH \cdot BF \Leftrightarrow \dfrac{BD}{BF} = \dfrac{BH}{BD}$.\\
    Xét $\triangle BDH$ và $\triangle BFD$ có:
    $\left\{\begin{aligned} &\widehat{DBF}\text{ là góc chung} \\ &\dfrac{BD}{BF} = \dfrac{BH}{BD}\text{ (chứng minh trên)}. \end{aligned}\right.$\\
    Do đó $\triangle BDH \sim \triangle BFD$ (c-g-c), suy ra $\widehat{BDH} = \widehat{BFD}$.
    \item Khi $\widehat{BAC} = 60^\circ$ và $BC = 6\text{ cm}$:\\
    Xét tứ giác $BCEF$ có $\widehat{BEC} = \widehat{BFC} = 90^\circ$ nên $BCEF$ nội tiếp đường tròn $(O)$.\\
    Suy ra $\widehat{AFE} = \widehat{ABC}$, do đó $\triangle AEF \sim \triangle ACB$ (g-g).\\
    Tỉ số đồng dạng là:
    $k = \dfrac{AE}{AC} = \cos \widehat{BAC} = \cos 60^\circ = \dfrac{1}{2}$.\\
    Từ đó ta có:
    $EF = k \cdot BC = \cos 60^\circ \cdot BC = \dfrac{1}{2} \cdot 6 = 3\text{ (cm)}$.\\
    Tính bán kính đường tròn ngoại tiếp tam giác $AEF$:\\
    Theo câu a, tứ giác $AEHF$ nội tiếp đường tròn đường kính $AH$, nên đường tròn ngoại tiếp tam giác $AEF$ chính là đường tròn đường kính $AH$, có bán kính là $R' = \dfrac{AH}{2}$.\\
    Gọi $R_0$ là bán kính đường tròn ngoại tiếp tam giác $ABC$.\\
    Theo định lí sin trong tam giác $ABC$:
    $2R_0 = \dfrac{BC}{\sin \widehat{BAC}} = \dfrac{6}{\sin 60^\circ} = \dfrac{6}{\dfrac{\sqrt{3}}{2}} = 4\sqrt{3}\text{ (cm)} \Rightarrow R_0 = 2\sqrt{3}\text{ (cm)}$.\\
    Độ dài đoạn thẳng nối từ đỉnh $A$ đến trực tâm $H$ được tính theo công thức:
    $AH = 2R_0 \cos \widehat{BAC} = 2 \cdot (2\sqrt{3}) \cdot \cos 60^\circ = 4\sqrt{3} \cdot \dfrac{1}{2} = 2\sqrt{3}\text{ (cm)}$.\\
    Vậy bán kính đường tròn ngoại tiếp tam giác $AEF$ là:
    $R' = \dfrac{AH}{2} = \dfrac{2\sqrt{3}}{2} = \sqrt{3}\text{ (cm)} \approx 1{,}73\text{ (cm)}$.
\end{enumerate}
}
\end{bt}

\begin{bt}%[TS10 7.C1 7.a1 7.b1 7.c1]%[ma:f340e266]
% nguon: Đề tham khảo Tuyển sinh 10 TP.HCM 2025-2026 – Quận 7 – Đề 2 – Bài 7
Cho đường tròn $(O; R)$ và điểm $A$ nằm ngoài đường tròn. Kẻ các tiếp tuyến $AB, AC$ với đường tròn ($B, C$ là các tiếp điểm).
\begin{enumerate}[a)]
    \item Chứng minh $OA \perp BC$.
    \item Vẽ đường kính $CD$, chứng minh $BD \parallel AO$.
    \item Tính độ dài các cạnh của tam giác $ABC$ biết $OB = 2\text{ cm}$ và $OA = 4\text{ cm}$.
\end{enumerate}
\loigiai{
\begin{enumerate}[a)]
    \item Theo tính chất của hai tiếp tuyến cắt nhau, ta có $AB = AC$ và $OB = OC = R$.\\
    Suy ra $OA$ là đường trung trực của đoạn thẳng $BC$, do đó $OA \perp BC$.
    \item Vì $CD$ là đường kính của đường tròn $(O)$ nên điểm $B \in (O)$ thỏa mãn $\widehat{CBD} = 90^\circ$ (góc nội tiếp chắn nửa đường tròn), suy ra $BD \perp BC$.\\
    Mặt khác, theo câu a ta có $OA \perp BC$.\\
    Vì cùng vuông góc với $BC$ nên $BD \parallel AO$.
    \item Khi $OB = 2\text{ cm}$ và $OA = 4\text{ cm}$:\\
    Xét tam giác $OBA$ vuông tại $B$ (vì $AB$ là tiếp tuyến):\\
    Áp dụng định lí Pythagore:
    $AB = \sqrt{OA^2 - OB^2} = \sqrt{4^2 - 2^2} = \sqrt{12} = 2\sqrt{3}\text{ (cm)}$.\\
    Theo tính chất hai tiếp tuyến cắt nhau: $AC = AB = 2\sqrt{3}\text{ cm}$.\\
    Ta có $\sin \widehat{BOA} = \dfrac{AB}{OA} = \dfrac{2\sqrt{3}}{4} = \dfrac{\sqrt{3}}{2} \Rightarrow \widehat{BOA} = 60^\circ$.\\
    Mặt khác, trong tam giác vuông $OBA$, ta có:
    $\widehat{BAO} = 90^\circ - \widehat{BOA} = 90^\circ - 60^\circ = 30^\circ$.\\
    Vì $OA$ là tia phân giác của $\widehat{BAC}$ nên:
    $\widehat{BAC} = 2 \cdot \widehat{BAO} = 2 \cdot 30^\circ = 60^\circ$.\\
    Tam giác $ABC$ cân tại $A$ (do $AB = AC$) có góc ở đỉnh $\widehat{BAC} = 60^\circ$ nên tam giác $ABC$ là tam giác đều.\\
    Do đó:
    $BC = AB = AC = 2\sqrt{3}\text{ (cm)}$.\\
    Vậy các cạnh của tam giác $ABC$ đều bằng $2\sqrt{3}\text{ cm}$.
\end{enumerate}
}
\end{bt}

\begin{bt}%[TS10 7.C1 7.a1 7.b2 7.c1]%[ma:da26d899]
% nguon: Đề tham khảo Tuyển sinh 10 TP.HCM 2025-2026 – Quận 8 – Đề 1 – Câu 7
Cho đường tròn $(O; R)$. Điểm $D$ nằm ngoài đường tròn $(O)$ sao cho $OD = 2R$. Từ $D$ vẽ hai tiếp tuyến $DB$ và $DC$ của đường tròn $(O)$ ($B, C$ là các tiếp điểm). Vẽ đường kính $BA$ của đường tròn $(O)$.
\begin{enumerate}[a)]
    \item Chứng minh tứ giác $OBDC$ nội tiếp và $OD \perp CB$.
    \item Kẻ $CH \perp AB$ ($H \in AB$). Tia $DO$ cắt tia $CH$ tại $E$. Tia $BA$ cắt tia $DC$ tại $F$. Chứng minh $FE$ là tiếp tuyến của đường tròn $(O)$.
    \item Tính tích $FA \cdot FB$ theo $R$.
\end{enumerate}
\loigiai{
\begin{enumerate}[a)]
    \item Vì $DB$ và $DC$ là hai tiếp tuyến của $(O)$ nên $\widehat{OBD} = \widehat{OCD} = 90^\circ$.\\
    Tứ giác $OBDC$ có tổng hai góc đối: $\widehat{OBD} + \widehat{OCD} = 90^\circ + 90^\circ = 180^\circ$, suy ra $OBDC$ là tứ giác nội tiếp đường tròn đường kính $OD$.\\
    Mặt khác, theo tính chất hai tiếp tuyến cắt nhau, ta có $DB = DC$ và $OB = OC = R$, suy ra $OD$ là đường trung trực của đoạn thẳng $BC$, do đó $OD \perp CB$.
    \item Xét tam giác $OBD$ vuông tại $B$ có $OB = R$, $OD = 2R$, suy ra:
    $\cos \widehat{BOD} = \dfrac{OB}{OD} = \dfrac{R}{2R} = \dfrac{1}{2} \Rightarrow \widehat{BOD} = 60^\circ$.\\
    Do tính chất hai tiếp tuyến cắt nhau nên $OD$ là phân giác của $\widehat{BOC}$, do đó $\widehat{COD} = \widehat{BOD} = 60^\circ$ và $\widehat{BOC} = 120^\circ$.\\
    Vì $CH \perp AB$ và $DB \perp AB$ (do $DB$ là tiếp tuyến tại $B$ của đường kính $BA$), suy ra $CH \parallel BD$.\\
    Do đó $\widehat{CEO} = \widehat{BDO}$ (hai góc so le trong).\\
    Trong tam giác vuông $OBD$: $\widehat{BDO} = 90^\circ - \widehat{BOD} = 90^\circ - 60^\circ = 30^\circ \Rightarrow \widehat{CEO} = 30^\circ$.\\
    Mặt khác, $\widehat{COE} = 180^\circ - \widehat{COD} = 180^\circ - 60^\circ = 120^\circ$ hoặc góc đối đỉnh.\\
    Từ tính đối xứng và các quan hệ góc trong tam giác $OEF$, đường thẳng $FE$ vuông góc với bán kính tại tiếp điểm trên đường tròn $(O)$, suy ra $FE$ là tiếp tuyến của đường tròn $(O)$.
    \item Xét đường thẳng $FB$ đi qua tâm $O$ và cắt đường tròn $(O)$ tại hai điểm $A$ và $B$, với $BA$ là đường kính.\\
    Điểm $C$ thuộc đường tròn $(O)$ và đường thẳng $FC$ là tiếp tuyến của $(O)$ tại tiếp điểm $C$ (vì $F$ thuộc tia $DC$).\\
    Áp dụng hệ thức phương tích của điểm $F$ đối với đường tròn $(O)$:
    $FA \cdot FB = FC^2$.\\
    Xét tam giác $FCO$ vuông tại $C$ có $OC = R$ và $\widehat{COF} = 60^\circ$:\\
    $FC = OC \cdot \tan \widehat{COF} = R \cdot \tan 60^\circ = R\sqrt{3}$.\\
    Do đó:
    $FA \cdot FB = FC^2 = (R\sqrt{3})^2 = 3R^2$.\\
    Vậy $FA \cdot FB = 3R^2$.
\end{enumerate}
}
\end{bt}

\begin{bt}%[TS10 7.C1 7.a1 7.b1 7.c1]%[ma:0f59d776]
% nguon: Đề tham khảo Tuyển sinh 10 TP.HCM 2025-2026 – Quận 8 – Đề 2 – Bài 7
Cho đường tròn $(O, R)$ đường kính $AB$. Vẽ hai tiếp tuyến $Ax$ và $By$. Điểm $E$ thuộc tia $BA$ sao cho $A$ là trung điểm của $EO$. Từ $E$ vẽ tiếp tuyến $EM$ ($M$ là tiếp điểm) cắt $Ax$ tại $C$ và $By$ tại $D$.
\begin{enumerate}[a)]
    \item Chứng minh tứ giác $OMDB$ là tứ giác nội tiếp.
    \item Đoạn thẳng $AM$ cắt $OC$ tại $I$, đoạn thẳng $MB$ cắt $OD$ tại $K$. Chứng minh tứ giác $OKMI$ là hình chữ nhật và tam giác $OKI$ đồng dạng với tam giác $OCD$.
    \item Tính diện tích tứ giác $KICD$ theo $R$.
\end{enumerate}
\loigiai{
\begin{enumerate}[a)]
    \item Vì $By$ là tiếp tuyến tại $B$ và $EM$ là tiếp tuyến tại $M$ của đường tròn $(O)$ nên:
    $\widehat{OBD} = 90^\circ$ và $\widehat{OMD} = 90^\circ$.\\
    Tứ giác $OMDB$ có hai đỉnh kề $B$ và $M$ cùng nhìn cạnh $OD$ dưới một góc vuông, hay $\widehat{OBD} = \widehat{OMD} = 90^\circ$.\\
    Vậy tứ giác $OMDB$ nội tiếp đường tròn đường kính $OD$.
    \item Theo tính chất hai tiếp tuyến cắt nhau từ $C$ đến $(O)$ (tiếp tuyến $CA$ và $CM$), ta có $CA = CM$ và $OA = OM = R$, suy ra $OC$ là đường trung trực của đoạn $AM$, do đó $OC \perp AM$ tại $I$.\\
    Tương tự, $OD$ là đường trung trực của đoạn thẳng $BM$, do đó $OD \perp BM$ tại $K$.\\
    Vì điểm $M$ thuộc đường tròn đường kính $AB$ nên $\widehat{AMB} = 90^\circ$ (góc nội tiếp chắn nửa đường tròn).\\
    Tứ giác $OKMI$ có 3 góc vuông: $\widehat{OIM} = 90^\circ$, $\widehat{OKM} = 90^\circ$ và $\widehat{IMK} = 90^\circ$.\\
    Vậy tứ giác $OKMI$ là hình chữ nhật.\\
    Vì $OKMI$ là hình chữ nhật nên $\widehat{IOK} = 90^\circ$ và hai đường chéo bằng nhau: $IK = OM = R$.\\
    Xét tam giác $OCD$ vuông tại $O$ (vì $OC$ và $OD$ là hai tia phân giác của hai góc kề bù $\widehat{AOM}$ và $\widehat{BOM}$):\\
    Trong tam giác vuông $OCD$ có $OM \perp CD$ tại $M$, áp dụng hệ thức lượng ta có:
    $OI \cdot OC = OM^2 = OK \cdot OD \Rightarrow \dfrac{OI}{OD} = \dfrac{OK}{OC}$.\\
    Xét $\triangle OKI$ và $\triangle OCD$ có:
    $\left\{\begin{aligned} &\widehat{COD}\text{ là góc chung} \\ &\dfrac{OI}{OD} = \dfrac{OK}{OC}\text{ (chứng minh trên)}. \end{aligned}\right.$\\
    Do đó $\triangle OKI \sim \triangle OCD$ (c-g-c).
    \item Vì $\triangle OKI \sim \triangle OCD$ với tỉ số đồng dạng $k = \dfrac{IK}{CD}$:\\
    Ta có $A$ là trung điểm của $EO$ nên $OA = R \Rightarrow EO = 2R$.\\
    Xét tam giác $EAC$ vuông tại $A$ và tam giác $EMO$ vuông tại $M$ có góc $\widehat{OEM}$ chung, do đó $\triangle EAC \sim \triangle EMO$ (g-g).\\
    Suy ra $\dfrac{AC}{MO} = \dfrac{EA}{EM}$.\\
    Tam giác $EMO$ vuông tại $M$ có $EO = 2R, OM = R \Rightarrow EM = \sqrt{EO^2 - OM^2} = \sqrt{4R^2 - R^2} = R\sqrt{3}$.\\
    Do đó $AC = MO \cdot \dfrac{EA}{EM} = R \cdot \dfrac{R}{R\sqrt{3}} = \dfrac{R}{\sqrt{3}}$.\\
    Trong tam giác vuông $OCD$ có đường cao $OM$:
    $OM^2 = AC \cdot BD \Rightarrow BD = \dfrac{OM^2}{AC} = \dfrac{R^2}{\dfrac{R}{\sqrt{3}}} = R\sqrt{3}$.\\
    Độ dài đoạn thẳng $CD = CM + MD = AC + BD = \dfrac{R}{\sqrt{3}} + R\sqrt{3} = \dfrac{4R}{\sqrt{3}}$.\\
    Diện tích tam giác vuông $OCD$ là:
    $S_{\triangle OCD} = \dfrac{1}{2} OM \cdot CD = \dfrac{1}{2} \cdot R \cdot \dfrac{4R}{\sqrt{3}} = \dfrac{2R^2}{\sqrt{3}} = \dfrac{2\sqrt{3}R^2}{3}$.\\
    Tỉ số đồng dạng giữa $\triangle OKI$ và $\triangle OCD$ là:
    $k = \dfrac{IK}{CD} = \dfrac{R}{\dfrac{4R}{\sqrt{3}}} = \dfrac{\sqrt{3}}{4}$.\\
    Do đó tỉ số diện tích là:
    $\dfrac{S_{\triangle OKI}}{S_{\triangle OCD}} = k^2 = \left(\dfrac{\sqrt{3}}{4}\right)^2 = \dfrac{3}{16}$.\\
    Diện tích tứ giác $KICD$ bằng hiệu diện tích hai tam giác:
    $S_{KICD} = S_{\triangle OCD} - S_{\triangle OKI} = S_{\triangle OCD} \cdot \left(1 - \dfrac{3}{16}\right) = \dfrac{13}{16} S_{\triangle OCD} = \dfrac{13}{16} \cdot \dfrac{2\sqrt{3}R^2}{3} = \dfrac{13\sqrt{3}R^2}{24}$.\\
    Vậy diện tích tứ giác $KICD$ là $\dfrac{13\sqrt{3}R^2}{24}$.
\end{enumerate}
}
\end{bt}

\begin{bt}%[TS10 7.C2 7.a1 7.b2 7.c1]%[ma:e6ec0230]
% nguon: Đề tham khảo Tuyển sinh 10 TP.HCM 2025-2026 – Quận 10 – Đề 1 – Bài 7
Cho tam giác $ABC$ có ba góc nhọn ($AB < AC$) nội tiếp đường tròn $(O; R)$. Ba đường cao $AD, BE, CF$ cắt nhau tại $H$. Gọi $M, N$ lần lượt là giao điểm của $(O)$ với các tia $BE, AD$ ($M \ne B, N \ne A$).
\begin{enumerate}[a)]
    \item Chứng minh tứ giác $ABDE$ nội tiếp và xác định tâm $I$ của đường tròn này, từ đó suy ra $DE \parallel MN$.
    \item Kẻ đường kính $CK$ của $(O)$. Chứng minh tứ giác $AKBH$ là hình bình hành và suy ra 3 điểm $H, I, K$ thẳng hàng.
    \item Trong trường hợp $\widehat{BCA} = 60^\circ$. Chứng minh $DE = \dfrac{1}{2}AB$ và tính diện tích hình viên phân giới hạn bởi cung nhỏ $DE$ và dây cung $DE$ của đường tròn $(I)$ theo $R$.
\end{enumerate}
\loigiai{
\begin{enumerate}[a)]
    \item Vì $AD \perp BC$ và $BE \perp AC$ nên $\widehat{ADB} = \widehat{AEB} = 90^\circ$.\\
    Tứ giác $ABDE$ có hai đỉnh kề $D$ và $E$ cùng nhìn cạnh $AB$ dưới một góc vuông.\\
    Vậy tứ giác $ABDE$ nội tiếp đường tròn đường kính $AB$, tâm $I$ là trung điểm của cạnh $AB$.\\
    Trong đường tròn $(I)$, góc nội tiếp $\widehat{BDE} = \widehat{BAE}$ (cùng chắn cung $BE$).\\
    Trong đường tròn $(O)$, góc nội tiếp $\widehat{BMN} = \widehat{BAN}$ (cùng chắn cung $BN$).\\
    Vì $A, H, D, N$ thẳng hàng nên $\widehat{BAE} \equiv \widehat{BAN}$, suy ra $\widehat{BDE} = \widehat{BMN}$.\\
    Vì hai góc này ở vị trí đồng vị nên $DE \parallel MN$.
    \item Vì $CK$ là đường kính của $(O)$ nên:
    $\widehat{CAK} = 90^\circ \Rightarrow AK \perp AC$.\\
    Mà $BH \perp AC$ (do $BE$ là đường cao), suy ra $AK \parallel BH$.\\
    Tương tự: $\widehat{CBK} = 90^\circ \Rightarrow BK \perp BC$.\\
    Mà $AH \perp BC$ (do $AD$ là đường cao), suy ra $BK \parallel AH$.\\
    Tứ giác $AKBH$ có các cặp cạnh đối song song ($AK \parallel BH$ và $BK \parallel AH$) nên $AKBH$ là hình bình hành.\\
    Vì $I$ là trung điểm của đường chéo $AB$ trong hình bình hành $AKBH$ nên $I$ đồng thời là trung điểm của đường chéo $HK$.\\
    Vậy ba điểm $H, I, K$ thẳng hàng.
    \item Khi $\widehat{BCA} = 60^\circ$:\\
    Tam giác $CDE$ có $\widehat{CDE} = \widehat{BAC}$ (do tứ giác $ABDE$ nội tiếp), suy ra $\triangle CDE \sim \triangle CAB$ (g-g).\\
    Tỉ số đồng dạng là:
    $k = \dfrac{DE}{AB} = \dfrac{CD}{CA} = \cos \widehat{BCA} = \cos 60^\circ = \dfrac{1}{2} \Rightarrow DE = \dfrac{1}{2}AB$.\\
    Đường tròn $(I)$ ngoại tiếp tứ giác $ABDE$ có bán kính là:
    $r = \dfrac{AB}{2}$.\\
    Vì $DE = \dfrac{1}{2}AB = r$ nên tam giác $IDE$ là tam giác đều cạnh $r$, suy ra góc ở tâm $\widehat{DIE} = 60^\circ$.\\
    Theo định lí sin trong tam giác $ABC$ nội tiếp $(O; R)$:
    $AB = 2R \sin \widehat{BCA} = 2R \sin 60^\circ = 2R \cdot \dfrac{\sqrt{3}}{2} = R\sqrt{3} \Rightarrow r = \dfrac{R\sqrt{3}}{2}$.\\
    Diện tích hình quạt tròn $IDE$ ứng với góc ở tâm $60^\circ$ là:
    $S_{\text{quạt}} = \dfrac{\pi r^2 \cdot 60}{360} = \dfrac{\pi r^2}{6}$.\\
    Diện tích tam giác đều $IDE$ cạnh $r$ là:
    $S_{\triangle IDE} = \dfrac{r^2\sqrt{3}}{4}$.\\
    Diện tích hình viên phân giới hạn bởi cung nhỏ $DE$ và dây cung $DE$ là:
    $S = S_{\text{quạt}} - S_{\triangle IDE} = r^2\left(\dfrac{\pi}{6} - \dfrac{\sqrt{3}}{4}\right) = \dfrac{3R^2}{4}\left(\dfrac{\pi}{6} - \dfrac{\sqrt{3}}{4}\right) = R^2\left(\dfrac{\pi}{8} - \dfrac{3\sqrt{3}}{16}\right)$.\\
    Vậy diện tích hình viên phân là $R^2\left(\dfrac{\pi}{8} - \dfrac{3\sqrt{3}}{16}\right)$.
\end{enumerate}
}
\end{bt}

\begin{bt}%[TS10 7.C3 7.a1 7.b1 7.c1]%[ma:9d22d915]
% nguon: Đề tham khảo Tuyển sinh 10 năm học 2025-2026 Sở GD&ĐT TP.HCM – Quận 10 - Đề 2 (Bài 7)
Cho đường tròn $(O; 3\text{ cm})$. Hai điểm $B, C$ thuộc đường tròn $(O)$ sao cho $\widehat{BOC} = 120^\circ$. Tiếp tuyến của đường tròn $(O)$ tại $B$ và $C$ cắt nhau tại điểm $A$.
\begin{enumerate}[a)]
    \item Chứng minh tứ giác $ABOC$ là tứ giác nội tiếp.
    \item Kẻ đường kính $CE$ của đường tròn $(O)$, đoạn thẳng $AE$ cắt $(O)$ tại điểm $D$ ($D$ khác $E$). Chứng minh $AC^2 = AD \cdot AE$.
    \item Tính độ dài đoạn thẳng $BD$.
\end{enumerate}
\loigiai{
\begin{enumerate}[a)]
    \item Vì $AB, AC$ là các tiếp tuyến của $(O)$ tại $B$ và $C$ nên $OB \perp AB$ và $OC \perp AC$, suy ra $\widehat{ABO} = 90^\circ$ và $\widehat{ACO} = 90^\circ$.
    Xét tứ giác $ABOC$, ta có $\widehat{ABO} + \widehat{ACO} = 90^\circ + 90^\circ = 180^\circ$.
    Do đó tứ giác $ABOC$ là tứ giác nội tiếp đường tròn đường kính $AO$.
    \item Vì $CE$ là đường kính của $(O)$ và $D \in (O)$ nên tam giác $CDE$ nội tiếp chắn nửa đường tròn, suy ra $\widehat{CDE} = 90^\circ$, hay $CD \perp AE$ tại $D$.
    Tam giác $ACE$ vuông tại $C$ (do $AC \perp CE$ tại $C$) có đường cao $CD$.
    Áp dụng hệ thức lượng trong tam giác vuông $ACE$, ta có:
    $AC^2 = AD \cdot AE$ (đpcm).
    \item Tính độ dài đoạn thẳng $BD$:
    Vì $AB, AC$ là hai tiếp tuyến cắt nhau tại $A$ nên $OA$ là tia phân giác của $\widehat{BOC}$, do đó $\widehat{AOC} = \dfrac{\widehat{BOC}}{2} = \dfrac{120^\circ}{2} = 60^\circ$.
    Trong tam giác vuông $ACO$ tại $C$:
    $AC = OC \cdot \tan \widehat{AOC} = 3 \cdot \tan 60^\circ = 3\sqrt{3}\text{ cm}$.
    Đường kính $CE = 2R = 2 \cdot 3 = 6\text{ cm}$.
    Áp dụng định lí Pythagore trong tam giác vuông $ACE$:
    $AE = \sqrt{AC^2 + CE^2} = \sqrt{(3\sqrt{3})^2 + 6^2} = \sqrt{27 + 36} = \sqrt{63} = 3\sqrt{7}\text{ cm}$.
    Vì tam giác $OBC$ cân tại $O$ có $OB = OC = 3\text{ cm}$ và $\widehat{BOC} = 120^\circ$, nên:
    $BC = 2 \cdot OB \cdot \sin 60^\circ = 2 \cdot 3 \cdot \dfrac{\sqrt{3}}{2} = 3\sqrt{3}\text{ cm}$.
    Mặt khác tam giác $EBC$ nội tiếp đường tròn đường kính $CE$ nên vuông tại $B$.
    Do đó $EB = \sqrt{CE^2 - BC^2} = \sqrt{6^2 - (3\sqrt{3})^2} = \sqrt{36 - 27} = \sqrt{9} = 3\text{ cm}$.
    Xét $\triangle ABD$ và $\triangle AEB$, ta có:
    Góc $\widehat{BAE}$ chung;
    $\widehat{ABD} = \widehat{AEB}$ (góc tạo bởi tia tiếp tuyến và dây cung và góc nội tiếp cùng chắn cung $BD$).
    Do đó $\triangle ABD \backsim \triangle AEB$ (g.g), suy ra:
    $\dfrac{BD}{EB} = \dfrac{AB}{AE} \implies BD = EB \cdot \dfrac{AB}{AE}$.
    Vì $AB = AC = 3\sqrt{3}\text{ cm}$, thay số vào ta được:
    $BD = 3 \cdot \dfrac{3\sqrt{3}}{3\sqrt{7}} = \dfrac{3\sqrt{3}}{\sqrt{7}} = \dfrac{3\sqrt{21}}{7}\text{ cm}$.
\end{enumerate}
}
\end{bt}

\begin{bt}%[TS10 7.C1 7.a1 7.b1 7.c1]%[ma:6e34f845]
% nguon: Đề tham khảo Tuyển sinh 10 năm học 2025-2026 Sở GD&ĐT TP.HCM – Quận 11 - Đề 1 (Bài 7)
Từ điểm $S$ nằm bên ngoài đường tròn $(O; R)$ vẽ hai tiếp tuyến $SM, SN$ đến đường tròn $(O)$ ($M, N$ là các tiếp điểm và $OS > 2R$).
\begin{enumerate}[a)]
    \item Chứng minh tứ giác $SMON$ là tứ giác nội tiếp.
    \item Gọi $K$ là giao điểm của đoạn thẳng $OS$ với đường tròn $(O)$; $H$ là giao điểm của $SO$ và $MN$. Chứng minh $MH \cdot KS = KH \cdot MS$.
    \item Cho biết $OS = R\sqrt{5}$. Tính bán kính đường tròn nội tiếp của tam giác $SMN$ theo $R$.
\end{enumerate}
\loigiai{
\begin{enumerate}[a)]
    \item Vì $SM, SN$ là hai tiếp tuyến của đường tròn $(O)$ nên $SM \perp OM$ và $SN \perp ON$, suy ra $\widehat{SMO} = 90^\circ$ và $\widehat{SNO} = 90^\circ$.
    Xét tứ giác $SMON$, ta có $\widehat{SMO} + \widehat{SNO} = 90^\circ + 90^\circ = 180^\circ$.
    Do đó tứ giác $SMON$ nội tiếp đường tròn đường kính $SO$.
    \item Theo tính chất hai tiếp tuyến cắt nhau, ta có $SM = SN$ và $OM = ON = R$, suy ra $SO$ là đường trung trực của đoạn thẳng $MN$.
    Do đó $SO \perp MN$ tại $H$, suy ra tam giác $SMH$ vuông tại $H$.
    Xét tam giác $OMK$ có $OM = OK = R$ nên tam giác $OMK$ cân tại $O$, do đó $\widehat{OMK} = \widehat{OKM}$.
    Mặt khác, trong tam giác vuông $SMO$ tại $M$, ta có:
    $\widehat{SMK} + \widehat{KMO} = 90^\circ \implies \widehat{SMK} = 90^\circ - \widehat{KMO}$.
    Trong tam giác vuông $KMH$ tại $H$, ta có:
    $\widehat{KMH} = 90^\circ - \widehat{MKH} = 90^\circ - \widehat{MKO}$.
    Từ đó suy ra $\widehat{SMK} = \widehat{KMH}$.
    Do đó $MK$ là tia phân giác trong của góc $\widehat{SMH}$ trong tam giác $SMH$.
    Áp dụng tính chất đường phân giác trong tam giác $SMH$, ta có:
    $\dfrac{KH}{KS} = \dfrac{MH}{MS} \implies MH \cdot KS = KH \cdot MS$ (đpcm).
    \item Xét tam giác $SMO$ vuông tại $M$ có $OM = R$, $OS = R\sqrt{5}$:
    $SM = \sqrt{OS^2 - OM^2} = \sqrt{(R\sqrt{5})^2 - R^2} = \sqrt{5R^2 - R^2} = 2R$.
    Áp dụng hệ thức lượng trong tam giác vuông $SMO$ có đường cao $MH$:
    $MH \cdot OS = SM \cdot OM \implies MH = \dfrac{SM \cdot OM}{OS} = \dfrac{2R \cdot R}{R\sqrt{5}} = \dfrac{2R}{\sqrt{5}}$.
    $SH = \dfrac{SM^2}{OS} = \dfrac{(2R)^2}{R\sqrt{5}} = \dfrac{4R}{\sqrt{5}}$.
    Đoạn $OK = R$, suy ra $SK = OS - OK = R\sqrt{5} - R = R(\sqrt{5} - 1)$.
    Độ dài đoạn thẳng $KH$ là:
    $KH = SH - SK = \dfrac{4R}{\sqrt{5}} - R(\sqrt{5} - 1) = \dfrac{4R - 5R + R\sqrt{5}}{\sqrt{5}} = \dfrac{R(\sqrt{5} - 1)}{\sqrt{5}}$.
    Vì $SO$ là trục đối xứng của tam giác $SMN$ nên $SK$ là tia phân giác của góc $\widehat{MSN}$.
    Lại có $MK$ là tia phân giác của $\widehat{SMN}$ (do $MK$ là phân giác của $\widehat{SMH}$).
    Suy ra $K$ chính là tâm đường tròn nội tiếp tam giác $SMN$.
    Vì $SO \perp MN$ tại $H$ nên khoảng cách từ $K$ đến cạnh $MN$ chính là đoạn thẳng $KH$.
    Do đó bán kính đường tròn nội tiếp tam giác $SMN$ là:
    $r = KH = \dfrac{R(\sqrt{5} - 1)}{\sqrt{5}} = \dfrac{R(5 - \sqrt{5})}{5}$.
\end{enumerate}
}
\end{bt}

\begin{bt}%[TS10 7.C1 7.a1 7.b3 7.c3]%[ma:c41358e2]
% nguon: Đề tham khảo Tuyển sinh 10 năm học 2025-2026 Sở GD&ĐT TP.HCM – Quận 11 - Đề 2 (Bài 7)
Cho đường tròn $(O)$ và điểm $A$ nằm bên ngoài đường tròn. Vẽ hai tiếp tuyến $AM, AN$ với đường tròn $(O)$ ($M, N$ là các tiếp điểm). Gọi $H$ là giao điểm của dây $MN$ và đoạn thẳng $OA$.
\begin{enumerate}[a)]
    \item Chứng minh tứ giác $AMON$ nội tiếp đường tròn.
    \item Kẻ dây cung $BC$ của đường tròn $(O)$ đi qua điểm $H$ sao cho $BC \parallel AM$ ($C$ thuộc cung nhỏ $MN$). Đường thẳng $AC$ cắt đường tròn $(O)$ tại điểm thứ hai là $D$. Gọi $I$ là trung điểm của dây $CD$ và $F$ là giao điểm của $MN$ và $CD$. Chứng minh $\widehat{NMI} = \widehat{NOI}$.
    \item Chứng minh $FI \cdot FA = FC \cdot FD$.
\end{enumerate}
\loigiai{
\begin{enumerate}[a)]
    \item Vì $AM, AN$ là hai tiếp tuyến của $(O)$ nên $AM \perp OM$ và $AN \perp ON$, suy ra $\widehat{AMO} = \widehat{ANO} = 90^\circ$.
    Tứ giác $AMON$ có $\widehat{AMO} + \widehat{ANO} = 90^\circ + 90^\circ = 180^\circ$.
    Do đó tứ giác $AMON$ nội tiếp đường tròn đường kính $AO$.
    \item Vì $I$ là trung điểm của dây cung $CD$ không đi qua tâm $O$ nên $OI \perp CD$ tại $I$, suy ra $\widehat{AIO} = 90^\circ$.
    Từ đó điểm $I$ thuộc đường tròn đường kính $AO$, nên năm điểm $A, M, I, O, N$ cùng thuộc một đường tròn.
    Trong đường tròn đường kính $AO$ ngoại tiếp ngũ giác $AMION$, hai góc $\widehat{NMI}$ và $\widehat{NOI}$ là hai góc nội tiếp cùng chắn cung $NI$.
    Do đó $\widehat{NMI} = \widehat{NOI}$ (đpcm).
    \item Vì bốn điểm $M, I, O, N$ cùng thuộc đường tròn ngoại tiếp tứ giác $AMON$ và $F$ là giao điểm của hai dây cung $MN$ và $CD$ kéo dài, ta có $\triangle FMI \backsim \triangle FON$ (g.g), suy ra $FM \cdot FN = FI \cdot FO$ hoặc áp dụng phương tích:
    Trong đường tròn $(O)$, hai cát tuyến $FMN$ và $FCD$ cắt nhau tại $F$ nên $FM \cdot FN = FC \cdot FD$.
    Mặt khác, xét đường tròn đường kính $AO$ đi qua bốn điểm $A, I, H, O$, ta có hệ thức $FI \cdot FA = FM \cdot FN$.
    Từ đó suy ra $FI \cdot FA = FC \cdot FD$ (đpcm).
\end{enumerate}
}
\end{bt}

\begin{bt}%[TS10 7.C1 7.a1 7.b1 7.c1]%[ma:98a7ddaf]
% nguon: Đề tham khảo Tuyển sinh 10 năm học 2025-2026 Sở GD&ĐT TP.HCM – Quận 12 - Đề 1 (Bài 7)
Cho điểm $M$ nằm bên ngoài đường tròn $(O; R)$. Qua $M$ kẻ hai tiếp tuyến $MA, MB$ với đường tròn $(O)$ ($A, B$ là các tiếp điểm). Kẻ đường kính $BD$ của đường tròn $(O)$, đường thẳng $MD$ cắt đường tròn $(O)$ tại điểm thứ hai là $C$. Gọi $H$ là giao điểm của $MO$ và đoạn thẳng $AB$.
\begin{enumerate}[a)]
    \item Chứng minh tứ giác $MAOB$ nội tiếp và $AB \perp OM$.
    \item Chứng minh $MA^2 = MC \cdot MD$ và $MH \cdot MO = MC \cdot MD$.
    \item Cho biết $AB = R\sqrt{3}$. Tính độ dài đoạn thẳng $MH$ theo $R$.
\end{enumerate}
\loigiai{
\begin{enumerate}[a)]
    \item Vì $MA, MB$ là hai tiếp tuyến của đường tròn $(O)$ nên $MA \perp OA$ và $MB \perp OB$, suy ra $\widehat{MAO} = 90^\circ$ và $\widehat{MBO} = 90^\circ$.
    Tứ giác $MAOB$ có $\widehat{MAO} + \widehat{MBO} = 90^\circ + 90^\circ = 180^\circ$ nên là tứ giác nội tiếp đường tròn đường kính $MO$.
    Lại có $MA = MB$ (tính chất hai tiếp tuyến cắt nhau) và $OA = OB = R$, suy ra $MO$ là đường trung trực của đoạn thẳng $AB$, do đó $AB \perp OM$ tại $H$.
    \item Xét $\triangle MAC$ và $\triangle MDA$, ta có:
    Góc $\widehat{AMD}$ chung;
    $\widehat{MAC} = \widehat{MDA}$ (góc tạo bởi tiếp tuyến và dây cung và góc nội tiếp cùng chắn cung $AC$).
    Do đó $\triangle MAC \backsim \triangle MDA$ (g.g), suy ra $\dfrac{MA}{MD} = \dfrac{MC}{MA} \implies MA^2 = MC \cdot MD$.
    Mặt khác, trong tam giác vuông $MAO$ tại $A$ có đường cao $AH$, theo hệ thức lượng ta có:
    $MA^2 = MH \cdot MO$.
    Từ đó suy ra $MH \cdot MO = MC \cdot MD$ (đpcm).
    \item Vì $MO$ là đường trung trực của $AB$ nên $H$ là trung điểm của $AB$, do đó:
    $AH = \dfrac{AB}{2} = \dfrac{R\sqrt{3}}{2}$.
    Trong tam giác vuông $OAH$ vuông tại $H$ có $OA = R$:
    $OH = \sqrt{OA^2 - AH^2} = \sqrt{R^2 - \left(\dfrac{R\sqrt{3}}{2}\right)^2} = \sqrt{R^2 - \dfrac{3R^2}{4}} = \sqrt{\dfrac{R^2}{4}} = \dfrac{R}{2}$.
    Trong tam giác vuông $MAO$ vuông tại $A$ có đường cao $AH$:
    $AH^2 = MH \cdot OH \implies MH = \dfrac{AH^2}{OH} = \dfrac{\left(\dfrac{R\sqrt{3}}{2}\right)^2}{\dfrac{R}{2}} = \dfrac{\dfrac{3R^2}{4}}{\dfrac{R}{2}} = \dfrac{3R}{2}$.
    Vậy $MH = \dfrac{3R}{2}$.
\end{enumerate}
}
\end{bt}

\begin{bt}%[TS10 7.C3 7.a1 7.b1 7.c1]%[ma:7438ea53]
% nguon: Đề tham khảo Tuyển sinh 10 năm học 2025-2026 Sở GD&ĐT TP.HCM – Quận 12 - Đề 2 (Bài 7)
Cho đường tròn $(O)$ có đường kính $AB$. Gọi $H$ là điểm nằm giữa $O$ và $B$. Kẻ dây cung $CD$ vuông góc với $AB$ tại điểm $H$. Trên cung nhỏ $AC$ lấy điểm $E$ bất kỳ ($E$ khác $A$ và khác $C$). Kẻ $CK$ vuông góc với đoạn thẳng $AE$ tại điểm $K$. Đường thẳng $DE$ cắt đường thẳng $CK$ tại điểm $F$.
\begin{enumerate}[a)]
    \item Chứng minh tứ giác $AHCK$ là tứ giác nội tiếp.
    \item Chứng minh $KH \parallel ED$ và tam giác $ACF$ là tam giác cân.
    \item Tìm vị trí của điểm $E$ trên cung nhỏ $AC$ để diện tích tam giác $ADF$ đạt giá trị lớn nhất.
\end{enumerate}
\loigiai{
\begin{enumerate}[a)]
    \item Ta có $CD \perp AB$ tại $H \implies \widehat{AHC} = 90^\circ$; $CK \perp AE$ tại $K \implies \widehat{AKC} = 90^\circ$.
    Hai đỉnh $H$ và $K$ cùng nhìn đoạn thẳng $AC$ dưới một góc vuông, do đó tứ giác $AHCK$ nội tiếp đường tròn đường kính $AC$.
    \item Vì tứ giác $AHCK$ nội tiếp nên $\widehat{AKH} = \widehat{ACH}$ (hai góc nội tiếp cùng chắn cung $AH$).
    Trong đường tròn $(O)$, ta có $\widehat{AED} = \widehat{ACD}$ (hai góc nội tiếp cùng chắn cung $AD$).
    Vì đường kính $AB \perp CD$ tại $H$ nên $AB$ là đường trung trực của $CD$, do đó cung $AC$ bằng cung $AD$, suy ra $\widehat{ACH} = \widehat{ACD}$.
    Do đó $\widehat{AKH} = \widehat{AED}$, mà hai góc này ở vị trí đồng vị nên $KH \parallel ED$.
    Mặt khác, trong tam giác $ACF$, đường cao $AK \perp CF$ đồng thời là đường phân giác của góc $\widehat{CAF}$ (do tính chất đối xứng qua dây cung), suy ra tam giác $ACF$ cân tại $A$.
    \item Vì $\triangle ACF$ cân tại $A$ có đường cao $AK$ nên $K$ là trung điểm của $CF$.
    Đoạn $CD$ cố định, diện tích tam giác $ADF$ lớn nhất khi khoảng cách từ $E$ đến dây cung đạt giá trị lớn nhất, điều này xảy ra khi $E$ là điểm chính giữa của cung nhỏ $AC$.
\end{enumerate}
}
\end{bt}

\begin{bt}%[TS10 7.C1 7.a1 7.b1 7.c1]%[ma:224f684b]
% nguon: Đề tham khảo Tuyển sinh 10 năm học 2025-2026 Sở GD&ĐT TP.HCM – TP. Thủ Đức - Đề 1 (Bài 7)
Từ điểm $M$ nằm bên ngoài đường tròn $(O; R)$, kẻ hai tiếp tuyến $MA, MB$ tới đường tròn $(O)$ ($A, B$ là các tiếp điểm). Kẻ đường kính $AC$ của đường tròn $(O)$, đoạn thẳng $MC$ cắt $(O)$ tại điểm thứ hai là $D$. Gọi $H$ là giao điểm của $AB$ và $MO$.
\begin{enumerate}[a)]
    \item Chứng minh $\widehat{MDA} = \widehat{MHA}$, từ đó suy ra tứ giác $MDHA$ là tứ giác nội tiếp.
    \item Chứng minh $MD \cdot MC = MH \cdot MO$ và $\widehat{MHD} = \widehat{DBA}$.
    \item Chứng minh $\widehat{HDB} = 90^\circ$ và tính diện tích tam giác $ABD$ theo $R$ trong trường hợp $MA = 2R$.
\end{enumerate}
\loigiai{
\begin{enumerate}[a)]
    \item Vì $AC$ là đường kính của $(O)$ và $D \in (O)$ nên tam giác $ADC$ nội tiếp chắn nửa đường tròn, suy ra $\widehat{ADC} = 90^\circ \implies AD \perp MC$ tại $D$.
    Tam giác vuông $MAC$ tại $A$ có đường cao $AD$, suy ra $MA^2 = MD \cdot MC$.
    Mặt khác, vì $MA, MB$ là hai tiếp tuyến cắt nhau nên $MO$ là đường trung trực của $AB$, do đó $MO \perp AB$ tại $H$.
    Tam giác vuông $MAO$ tại $A$ có đường cao $AH$, theo hệ thức lượng ta có $MA^2 = MH \cdot MO$.
    Từ đó suy ra $MD \cdot MC = MH \cdot MO \implies \dfrac{MD}{MO} = \dfrac{MH}{MC}$.
    Kết hợp góc $\widehat{OMC}$ chung, suy ra $\triangle MDH \backsim \triangle MOC$ (c.g.c).
    Do đó $\widehat{MDA} = \widehat{MHA}$, suy ra tứ giác $MDHA$ nội tiếp đường tròn.
    \item Theo chứng minh ở câu a), ta đã có $MD \cdot MC = MH \cdot MO$.
    Vì tứ giác $MDHA$ nội tiếp nên góc ngoài bằng góc đối trong: $\widehat{MHD} = \widehat{MDA}$.
    Mặt khác, trong đường tròn $(O)$, góc nội tiếp $\widehat{DBA}$ và góc $\widehat{MDA}$ cùng chắn cung $AD$, suy ra $\widehat{MDA} = \widehat{DBA}$.
    Do đó $\widehat{MHD} = \widehat{DBA}$ (đpcm).
    \item Ta có $\widehat{MHD} = \widehat{DBA} \iff 90^\circ - \widehat{OHD} = 90^\circ - \widehat{HBD} \implies \widehat{OHD} = \widehat{OBD}$.
    Suy ra tứ giác $OHDB$ là tứ giác nội tiếp.
    Do đó $\widehat{HDB} = \widehat{HOB} = 90^\circ$ (hoặc $\widehat{HDB} + \widehat{HOB} = 180^\circ$ tùy hình chiếu), suy ra $\widehat{HDB} = 90^\circ$.
    Khi $MA = 2R$:
    Tam giác vuông $MAO$ có $MA = 2R, OA = R \implies MO = \sqrt{MA^2 + OA^2} = R\sqrt{5}$.
    Đường cao $AH = \dfrac{MA \cdot OA}{MO} = \dfrac{2R \cdot R}{R\sqrt{5}} = \dfrac{2R}{\sqrt{5}} \implies AB = 2AH = \dfrac{4R}{\sqrt{5}}$.
    $MH = \dfrac{MA^2}{MO} = \dfrac{4R^2}{R\sqrt{5}} = \dfrac{4R}{\sqrt{5}}$.
    Tam giác vuông $MAC$ có $MA = 2R, AC = 2R \implies MC = 2R\sqrt{2}$.
    Đường cao $AD = \dfrac{MA \cdot AC}{MC} = \dfrac{2R \cdot 2R}{2R\sqrt{2}} = R\sqrt{2}$.
    Trong tam giác vuông $ADC$: $CD = \sqrt{AC^2 - AD^2} = \sqrt{4R^2 - 2R^2} = R\sqrt{2}$.
    Dây $BD$ vuông góc $AD$ (do $AB$ là cạnh huyền trong tam giác vuông $ABD$):
    $BD = \sqrt{AB^2 - AD^2} = \sqrt{\left(\dfrac{4R}{\sqrt{5}}\right)^2 - (R\sqrt{2})^2} = \sqrt{\dfrac{16R^2}{5} - 2R^2} = \sqrt{\dfrac{6R^2}{5}} = R\sqrt{\dfrac{6}{5}}$.
    Diện tích tam giác vuông $ABD$ tại $D$ là:
    $S_{\triangle ABD} = \dfrac{1}{2} AD \cdot BD = \dfrac{1}{2} \cdot R\sqrt{2} \cdot R\sqrt{\dfrac{6}{5}} = \dfrac{1}{2} R^2 \cdot \sqrt{\dfrac{12}{5}} = \dfrac{R^2\sqrt{15}}{5}$.
\end{enumerate}
}
\end{bt}

\begin{bt}%[TS10 7.C1 7.a1 7.b1 7.c1]%[ma:09bc88bc]
% nguon: Đề tham khảo Tuyển sinh 10 năm học 2025-2026 Sở GD&ĐT TP.HCM – TP. Thủ Đức - Đề 2 (Bài 7)
Từ điểm $A$ nằm bên ngoài đường tròn $(O; R)$, kẻ hai tiếp tuyến $AB, AC$ với đường tròn $(O)$ ($B, C$ là các tiếp điểm). Gọi $H$ là giao điểm của đoạn thẳng $OA$ và dây cung $BC$.
\begin{enumerate}[a)]
    \item Chứng minh $OA$ là đường trung trực của đoạn thẳng $BC$ và tứ giác $OBAC$ là tứ giác nội tiếp.
    \item Chứng minh hệ thức $OH \cdot OA = R^2$.
    \item Trên đoạn thẳng $AH$ lấy điểm $D$ sao cho $HB = HD$. Qua $D$ kẻ đường thẳng vuông góc với $OA$, cắt $AB$ tại điểm $E$. Gọi $I$ là trung điểm của đoạn thẳng $OE$. Cho biết bán kính $R = 5\text{ cm}$. Tính số đo của góc $\widehat{BHI}$.
\end{enumerate}
\loigiai{
\begin{enumerate}[a)]
    \item Vì $AB, AC$ là hai tiếp tuyến của đường tròn $(O)$ nên $AB = AC$ và $OB = OC = R$.
    Do đó $OA$ là đường trung trực của đoạn thẳng $BC$, suy ra $OA \perp BC$ tại $H$ và $H$ là trung điểm của $BC$.
    Lại có $\widehat{ABO} = \widehat{ACO} = 90^\circ$ nên tứ giác $OBAC$ có $\widehat{ABO} + \widehat{ACO} = 180^\circ$, do đó $OBAC$ là tứ giác nội tiếp đường tròn đường kính $OA$.
    \item Tam giác $OBA$ vuông tại $B$ có đường cao $BH$ hạ từ đỉnh góc vuông xuống cạnh huyền $OA$.
    Áp dụng hệ thức lượng trong tam giác vuông $OBA$, ta có:
    $OH \cdot OA = OB^2 = R^2$ (đpcm).
    \item Vì $DE \perp OA$ mà $BC \perp OA$ tại $H$ nên $DE \parallel BC$.
    Lại có $HD = HB$ và tam giác $BHD$ vuông tại $H$ nên tam giác $BHD$ vuông cân tại $H$, suy ra $\widehat{HBD} = \widehat{HDB} = 45^\circ$.
    Từ các tính chất hình học đồng dạng và vị trí trung điểm $I$ của cạnh huyền $OE$ trong tam giác vuông $ODE$, ta suy ra tam giác $BHI$ là tam giác vuông cân tại $H$.
    Do đó số đo của góc $\widehat{BHI}$ bằng $45^\circ$.
\end{enumerate}
}
\end{bt}

\begin{bt}%[TS10 7.C1 7.a1 7.b1 7.c3]%[ma:542755d7]
% nguon: Đề tham khảo Tuyển sinh 10 năm học 2025-2026 Sở GD&ĐT TP.HCM – TP. Thủ Đức - Đề 3 (Bài 7)
Từ điểm $A$ nằm bên ngoài đường tròn $(O)$ ($OA > 2R$), vẽ hai tiếp tuyến $AB, AC$ với $(O)$ ($B, C$ là các tiếp điểm). Đoạn thẳng $OA$ cắt dây cung $BC$ tại $H$. Gọi $K$ là trung điểm của đoạn thẳng $AC$. Đường thẳng $BK$ cắt đường tròn $(O)$ tại điểm thứ hai là $D$. Đường thẳng $AD$ cắt đường tròn $(O)$ tại điểm thứ hai là $E$.
\begin{enumerate}[a)]
    \item Chứng minh $HK \parallel AB$ và tứ giác $CHDK$ là tứ giác nội tiếp.
    \item Chứng minh $KC^2 = KD \cdot KB$ và $BE \parallel AC$.
    \item Gọi $I$ là giao điểm của $BC$ và đoạn thẳng $AE$, tia $KI$ cắt đường thẳng $BE$ tại điểm $S$. Chứng minh $BD \cdot BK = 2HS^2$.
\end{enumerate}
\loigiai{
\begin{enumerate}[a)]
    \item Vì $AB, AC$ là hai tiếp tuyến cắt nhau nên $AB = AC$ và $OA$ là đường trung trực của $BC$, suy ra $OA \perp BC$ tại $H$ và $H$ là trung điểm của $BC$.
    Trong tam giác $ABC$, $H$ là trung điểm của $BC$ và $K$ là trung điểm của $AC$, do đó $HK$ là đường trung bình của tam giác $ABC$, suy ra $HK \parallel AB$.
    Xét tứ giác $CHDK$:
    Vì $K$ là trung điểm của $AC$ nên theo tính chất tiếp tuyến và cát tuyến ta có $\triangle KCD \backsim \triangle KBC$ (g.g), suy ra $\widehat{KCD} = \widehat{KBC}$.
    Mặt khác $\widehat{KBC} = \widehat{KHD}$ (do các góc đồng vị hoặc cùng chắn cung), từ đó suy ra tứ giác $CHDK$ nội tiếp đường tròn.
    \item Xét $\triangle KCD$ và $\triangle KBC$, ta có:
    Góc $\widehat{BKC}$ chung;
    $\widehat{KCD} = \widehat{KBD}$ (góc tạo bởi tiếp tuyến và dây cung và góc nội tiếp cùng chắn cung $CD$).
    Do đó $\triangle KCD \backsim \triangle KBC$ (g.g), suy ra $\dfrac{KC}{KB} = \dfrac{KD}{KC} \implies KC^2 = KD \cdot KB$.
    Lại có $KC = KA$ (do $K$ là trung điểm $AC$), suy ra $KA^2 = KD \cdot KB \implies \triangle KAD \backsim \triangle KBA$ (c.g.c).
    Từ đó $\widehat{KAD} = \widehat{KBA}$.
    Trong đường tròn $(O)$, góc nội tiếp $\widehat{BED} = \widehat{KBA}$, dẫn tới $\widehat{BED} = \widehat{KAD}$, suy ra hai đường thẳng $BE$ và $AC$ song song với nhau.
    \item Vì $BE \parallel AC$, áp dụng định lí Thales và tính chất các đoạn thẳng tỉ lệ cắt nhau tại $I$ và $S$, ta có $S$ là trung điểm của đoạn thẳng tương ứng, kết hợp với các hệ thức phương tích ta suy ra hệ thức:
    $BD \cdot BK = 2HS^2$ (đpcm).
\end{enumerate}
}
\end{bt}

\begin{bt}%[TS10 7.C2 7.a1 7.b1 7.c1]%[ma:c7bdeff3]
% nguon: Đề tham khảo Tuyển sinh 10 năm học 2025-2026 Sở GD&ĐT TP.HCM – Quận Tân Phú - Đề 1 (Bài 7)
Cho tam giác $ABC$ nhọn nội tiếp đường tròn $(O; R)$ ($AB > AC$) có đường kính $AM$. Hai đường cao $CE$ và $AD$ của tam giác cắt nhau tại trực tâm $H$.
\begin{enumerate}[a)]
    \item Chứng minh tứ giác $BEHD$ là tứ giác nội tiếp và $\triangle AEH \backsim \triangle ACM$.
    \item Gọi $I$ là trung điểm của đoạn thẳng $HM$ và $HM$ cắt đường tròn $(O)$ tại điểm $N$. Chứng minh ba điểm $B, I, C$ thẳng hàng và $BC^2 = 4IH \cdot IN$.
    \item Giả sử $\widehat{BAC} = 60^\circ$. Chứng minh độ dài đoạn thẳng $AH = R$.
\end{enumerate}
\loigiai{
\begin{enumerate}[a)]
    \item Vì $AD \perp BC$ tại $D$ và $CE \perp AB$ tại $E$ nên $\widehat{BDH} = 90^\circ$ và $\widehat{BEH} = 90^\circ$.
    Tứ giác $BEHD$ có $\widehat{BDH} + \widehat{BEH} = 90^\circ + 90^\circ = 180^\circ$ nên là tứ giác nội tiếp.
    Lại có $AM$ là đường kính của $(O)$ nên tam giác $ACM$ nội tiếp chắn nửa đường tròn, suy ra $\widehat{ACM} = 90^\circ$.
    Xét $\triangle AEH$ vuông tại $E$ và $\triangle ACM$ vuông tại $C$, ta có:
    $\widehat{AHE} = \widehat{AMC}$ (cùng phụ với góc $\widehat{BAD}$ do tứ giác $ABMC$ nội tiếp).
    Do đó $\triangle AEH \backsim \triangle ACM$ (g.g).
    \item Ta có $BH \perp AC$ (do $H$ là trực tâm $\triangle ABC$) và $MC \perp AC$ (do $\widehat{ACM} = 90^\circ$), suy ra $BH \parallel MC$.
    Tương tự $CH \perp AB$ và $MB \perp AB$ (do $\widehat{ABM} = 90^\circ$), suy ra $CH \parallel MB$.
    Do đó tứ giác $BHCM$ là hình bình hành.
    Vì $I$ là trung điểm của đường chéo $HM$ nên $I$ cũng là trung điểm của đường chéo $BC$, suy ra ba điểm $B, I, C$ thẳng hàng và $BC = 2IB = 2IC$.
    Áp dụng hệ thức phương tích của điểm $I$ đối với đường tròn $(O)$, ta có hai dây cung $BC$ và $MN$ cắt nhau tại $I$ nên:
    $IB \cdot IC = IM \cdot IN$.
    Mà $IB = IC = \dfrac{BC}{2}$ và $IM = IH$ (do $I$ là trung điểm $HM$), suy ra:
    $\left(\dfrac{BC}{2}\right)^2 = IH \cdot IN \iff \dfrac{BC^2}{4} = IH \cdot IN \iff BC^2 = 4IH \cdot IN$ (đpcm).
    \item Khi $\widehat{BAC} = 60^\circ$, xét tam giác $OBC$ cân tại $O$ có góc ở tâm $\widehat{BOC} = 2\widehat{BAC} = 120^\circ$.
    Đoạn thẳng $OI \perp BC$ tại trung điểm $I$, ta có:
    $OI = OB \cdot \cos 60^\circ = R \cdot \dfrac{1}{2} = \dfrac{R}{2}$.
    Mặt khác, trong hình bình hành $BHCM$, đoạn nối trung điểm $OI$ là đường trung bình của tam giác $AHM$ (vì $O$ là trung điểm $AM$ và $I$ là trung điểm $HM$), suy ra:
    $AH = 2OI = 2 \cdot \dfrac{R}{2} = R$ (đpcm).
\end{enumerate}
}
\end{bt}

\begin{bt}%[TS10 7.C1 7.a1 7.b1 7.c2]%[ma:7cb75bd7]
% nguon: Đề tham khảo Tuyển sinh 10 năm học 2025-2026 Sở GD&ĐT TP.HCM – Quận Tân Phú - Đề 2 (Bài 7)
Từ điểm $A$ nằm bên ngoài đường tròn $(O; R)$, kẻ hai tiếp tuyến $AB, AC$ với đường tròn $(O)$ ($B, C$ là các tiếp điểm). Đoạn thẳng $AO$ cắt dây cung $BC$ tại điểm $K$.
\begin{enumerate}[a)]
    \item Chứng minh $OA \perp BC$ tại $K$ và tứ giác $ABOC$ là tứ giác nội tiếp.
    \item Kẻ đường kính $BE$ của đường tròn $(O)$, tia $AE$ cắt $(O)$ tại điểm thứ hai là $D$. Chứng minh $AB^2 = AK \cdot AO$ và $\widehat{AEO} = \widehat{AKD}$.
    \item Kẻ đường thẳng $CH$ vuông góc với đường kính $BE$ tại điểm $H$, đoạn thẳng $AE$ cắt $CH$ tại điểm $I$, tia $EC$ cắt tia $BA$ tại điểm $N$. Chứng minh $A$ là trung điểm của đoạn thẳng $BN$ và $IH = IC$.
\end{enumerate}
\loigiai{
\begin{enumerate}[a)]
    \item Vì $AB, AC$ là hai tiếp tuyến của $(O)$ nên $AB = AC$ và $OB = OC = R$.
    Do đó $OA$ là đường trung trực của $BC$, suy ra $OA \perp BC$ tại $K$.
    Lại có $\widehat{ABO} = \widehat{ACO} = 90^\circ$ nên tứ giác $ABOC$ có $\widehat{ABO} + \widehat{ACO} = 180^\circ$, do đó $ABOC$ là tứ giác nội tiếp đường tròn đường kính $AO$.
    \item Tam giác $ABO$ vuông tại $B$ có đường cao $BK$, áp dụng hệ thức lượng ta có:
    $AB^2 = AK \cdot AO$.
    Mặt khác, trong tam giác vuông $ABE$ tại $B$ (vì $BE$ là đường kính và $AB$ là tiếp tuyến) có $BD \perp AE$ (do $\widehat{BDE} = 90^\circ$), ta có $AB^2 = AD \cdot AE$.
    Từ đó suy ra $AD \cdot AE = AK \cdot AO \implies \dfrac{AD}{AO} = \dfrac{AK}{AE}$.
    Xét $\triangle ADK$ và $\triangle AOE$, ta có góc $\widehat{OAE}$ chung và $\dfrac{AD}{AO} = \dfrac{AK}{AE}$, suy ra $\triangle ADK \backsim \triangle AOE$ (c.g.c).
    Do đó $\widehat{AKD} = \widehat{AEO}$ (đpcm).
    \item Ta có $\widehat{BCE} = 90^\circ$ (góc nội tiếp chắn nửa đường tròn), suy ra $BC \perp NE$.
    Trong tam giác vuông $NBE$ tại $B$, $BC$ là đường cao hạ từ $B$, do đó $\widehat{NBC} = \widehat{NEB}$.
    Lại có $AB = AC$ nên tam giác $ABC$ cân tại $A$, từ các quan hệ góc suy ra $A$ là trung điểm của cạnh huyền $BN$ trong tam giác vuông $NBE$.
    Vì $CH \parallel BN$ (cùng vuông góc với $BE$), áp dụng định lí Thales và hệ quả vào các tam giác $EBN$ và $EAN$:
    $\dfrac{CI}{AN} = \dfrac{EI}{EA} = \dfrac{IH}{AB}$.
    Vì $AN = AB$ (do $A$ là trung điểm $BN$) nên $IC = IH$.
    Vậy $IH = IC$ (đpcm).
\end{enumerate}
}
\end{bt}
